\documentclass[11pt]{article}

% ---- theme variables (passed via pandoc -V) ----
\newcommand{\ThemeName}{Deep Learning}
\newcommand{\ThemeKicker}{AI/ML Track}
\newcommand{\ThemeNumber}{03}

\usepackage{interviewstyle}
\setthemecolor{7c3aed}{a78bfa}

% pandoc-required plumbing
\providecommand{\tightlist}{\setlength{\itemsep}{1pt}\setlength{\parskip}{0pt}}
\setlength{\emergencystretch}{3em}
\providecommand{\pandocbounded}[1]{#1}

% syntax-highlighting macros emitted by pandoc (skylighting)
\usepackage{color}
\usepackage{fancyvrb}
\newcommand{\VerbBar}{|}
\newcommand{\VERB}{\Verb[commandchars=\\\{\}]}
\DefineVerbatimEnvironment{Highlighting}{Verbatim}{commandchars=\\\{\}}
% Add ',fontsize=\small' for more characters per line
\usepackage{framed}
\definecolor{shadecolor}{RGB}{35,38,41}
\newenvironment{Shaded}{\begin{snugshade}}{\end{snugshade}}
\newcommand{\AlertTok}[1]{\textcolor[rgb]{0.58,0.85,0.30}{\textbf{\colorbox[rgb]{0.30,0.12,0.14}{#1}}}}
\newcommand{\AnnotationTok}[1]{\textcolor[rgb]{0.25,0.50,0.35}{#1}}
\newcommand{\AttributeTok}[1]{\textcolor[rgb]{0.16,0.50,0.73}{#1}}
\newcommand{\BaseNTok}[1]{\textcolor[rgb]{0.96,0.45,0.00}{#1}}
\newcommand{\BuiltInTok}[1]{\textcolor[rgb]{0.50,0.55,0.55}{#1}}
\newcommand{\CharTok}[1]{\textcolor[rgb]{0.24,0.68,0.91}{#1}}
\newcommand{\CommentTok}[1]{\textcolor[rgb]{0.48,0.49,0.49}{#1}}
\newcommand{\CommentVarTok}[1]{\textcolor[rgb]{0.50,0.55,0.55}{#1}}
\newcommand{\ConstantTok}[1]{\textcolor[rgb]{0.15,0.68,0.68}{\textbf{#1}}}
\newcommand{\ControlFlowTok}[1]{\textcolor[rgb]{0.99,0.74,0.29}{\textbf{#1}}}
\newcommand{\DataTypeTok}[1]{\textcolor[rgb]{0.16,0.50,0.73}{#1}}
\newcommand{\DecValTok}[1]{\textcolor[rgb]{0.96,0.45,0.00}{#1}}
\newcommand{\DocumentationTok}[1]{\textcolor[rgb]{0.64,0.20,0.25}{#1}}
\newcommand{\ErrorTok}[1]{\textcolor[rgb]{0.85,0.27,0.33}{\underline{#1}}}
\newcommand{\ExtensionTok}[1]{\textcolor[rgb]{0.00,0.60,1.00}{\textbf{#1}}}
\newcommand{\FloatTok}[1]{\textcolor[rgb]{0.96,0.45,0.00}{#1}}
\newcommand{\FunctionTok}[1]{\textcolor[rgb]{0.56,0.27,0.68}{#1}}
\newcommand{\ImportTok}[1]{\textcolor[rgb]{0.15,0.68,0.38}{#1}}
\newcommand{\InformationTok}[1]{\textcolor[rgb]{0.77,0.36,0.00}{#1}}
\newcommand{\KeywordTok}[1]{\textcolor[rgb]{0.81,0.81,0.76}{\textbf{#1}}}
\newcommand{\NormalTok}[1]{\textcolor[rgb]{0.81,0.81,0.76}{#1}}
\newcommand{\OperatorTok}[1]{\textcolor[rgb]{0.81,0.81,0.76}{#1}}
\newcommand{\OtherTok}[1]{\textcolor[rgb]{0.15,0.68,0.38}{#1}}
\newcommand{\PreprocessorTok}[1]{\textcolor[rgb]{0.15,0.68,0.38}{#1}}
\newcommand{\RegionMarkerTok}[1]{\textcolor[rgb]{0.16,0.50,0.73}{\colorbox[rgb]{0.08,0.19,0.26}{#1}}}
\newcommand{\SpecialCharTok}[1]{\textcolor[rgb]{0.24,0.68,0.91}{#1}}
\newcommand{\SpecialStringTok}[1]{\textcolor[rgb]{0.85,0.27,0.33}{#1}}
\newcommand{\StringTok}[1]{\textcolor[rgb]{0.96,0.31,0.31}{#1}}
\newcommand{\VariableTok}[1]{\textcolor[rgb]{0.15,0.68,0.68}{#1}}
\newcommand{\VerbatimStringTok}[1]{\textcolor[rgb]{0.85,0.27,0.33}{#1}}
\newcommand{\WarningTok}[1]{\textcolor[rgb]{0.85,0.27,0.33}{#1}}

\begin{document}

\makecoverpage

\thispagestyle{empty}
{\hypersetup{hidelinks}\tableofcontents}
\clearpage

\begin{quote}
A neural network is a universal function approximator built from stacked, differentiable transformations --- trained end-to-end by propagating error gradients backward through the chain rule.
\end{quote}

\section{Overview --- What It Is}\label{overview--what-it-is}

Deep Learning is a subfield of machine learning that uses \textbf{multi-layered artificial neural networks} to learn hierarchical representations directly from raw data. Instead of handcrafting features, the network learns them automatically through a training process that adjusts millions (or billions) of parameters.

\textbf{Key idea:} Stack many parametric transformations (layers), apply non-linear activation functions between them, define a loss measuring prediction quality, and iteratively minimize the loss via gradient descent by flowing gradients backward through every layer.

\textbf{"Deep"} = more than one hidden layer. Depth enables the network to compose simple features into complex ones: edges \(\rightarrow\) textures \(\rightarrow\) parts \(\rightarrow\) objects (in vision); characters \(\rightarrow\) morphemes \(\rightarrow\) words \(\rightarrow\) meaning (in NLP).

\section{Why It Exists}\label{why-it-exists}

\begin{longtable}[]{@{}
  >{\raggedright\arraybackslash}p{(\columnwidth - 4\tabcolsep) * \real{0.1221}}
  >{\raggedright\arraybackslash}p{(\columnwidth - 4\tabcolsep) * \real{0.4390}}
  >{\raggedright\arraybackslash}p{(\columnwidth - 4\tabcolsep) * \real{0.4390}}@{}}
\toprule\noalign{}
\rowcolor{acc!16}
\begin{minipage}[b]{\linewidth}\raggedright
Era
\end{minipage} & \begin{minipage}[b]{\linewidth}\raggedright
Limitation
\end{minipage} & \begin{minipage}[b]{\linewidth}\raggedright
Deep Learning Solution
\end{minipage} \\
\midrule\noalign{}
\endhead
\bottomrule\noalign{}
\endlastfoot
Pre-1980s & Linear models can\textquotesingle t capture non-linear patterns & Non-linear activations + multiple layers \\
\rowcolor{zebra}
1980s--2000s & Shallow nets, vanishing gradients, limited data & Better init, ReLU, more data \\
2006--2012 & Hand-engineered features dominate & Pre-training, GPUs, AlexNet (2012) breakthrough \\
\rowcolor{zebra}
2012--present & Dataset scale + compute & Large-scale training on GPUs/TPUs, transformers \\
\end{longtable}

Deep Learning exists because:

\begin{enumerate}
\def\labelenumi{\arabic{enumi}.}
\tightlist
\item
  \textbf{Raw data is high-dimensional} --- pixels, tokens, audio samples. Hand-engineered features are brittle.
\item
  \textbf{Universal approximation theorem} --- an MLP with one hidden layer can approximate any continuous function. Depth makes this efficient.
\item
  \textbf{GPU parallelism} --- matrix multiplications are embarrassingly parallel.
\item
  \textbf{Data availability} --- ImageNet (1.2M images), web text corpora, etc.
\end{enumerate}

\section{Why FAANG Cares (Be Specific)}\label{why-faang-cares-be-specific}

\begin{itemize}
\tightlist
\item
  \textbf{Google/DeepMind} --- Search ranking, Google Translate (Transformer), AlphaFold, Gemini. Every ranking model at Google is a deep neural net.
\item
  \textbf{Meta/Facebook} --- Recommendation systems (DLRM), content moderation (image/video classifiers), ads CTR models, Llama LLMs.
\item
  \textbf{Amazon} --- Alexa (RNN/Transformer ASR), product recommendation, fraud detection, Alexa voice embeddings.
\item
  \textbf{Apple} --- Face ID (CNN embeddings), Siri (seq2seq), on-device ML (CoreML), photo classification.
\item
  \textbf{Netflix} --- Recommendation via deep collaborative filtering, thumbnail A/B test click-through CNNs.
\item
  \textbf{OpenAI/Microsoft} --- GPT models, Copilot, Azure ML infrastructure.
\end{itemize}

\textbf{Why interviewers probe this deeply:}

\begin{itemize}
\tightlist
\item
  Deep learning engineers need to debug training (exploding/vanishing gradients, overfitting, dead neurons).
\item
  Architectural choices (CNN vs RNN vs Transformer) have large cost/quality trade-offs at FAANG scale.
\item
  Knowing backprop from first principles = knowing how to debug gradient flow issues.
\item
  Knowing normalization + initialization = avoiding 90\% of training instability issues.
\end{itemize}

\section{Core Concepts}\label{core-concepts}

\subsection{The Perceptron --- Atomic Unit}\label{the-perceptron--atomic-unit}

The \textbf{perceptron} (Rosenblatt, 1958) is a single artificial neuron:

\setcodelang{Python}

\begin{Shaded}
\begin{Highlighting}[]
\NormalTok{y }\OperatorTok{=}\NormalTok{ f( w1}\OperatorTok{*}\NormalTok{x1 }\OperatorTok{+}\NormalTok{ w2}\OperatorTok{*}\NormalTok{x2 }\OperatorTok{+}\NormalTok{ ... }\OperatorTok{+}\NormalTok{ wn}\OperatorTok{*}\NormalTok{xn }\OperatorTok{+}\NormalTok{ b )}
  \OperatorTok{=}\NormalTok{ f( W·x }\OperatorTok{+}\NormalTok{ b )}
\end{Highlighting}
\end{Shaded}

\begin{itemize}
\tightlist
\item
  \textbf{x} = input vector
\item
  \textbf{W} = weight vector (what the network learns)
\item
  \textbf{b} = bias scalar (shifts the decision boundary)
\item
  \textbf{f} = activation function (adds non-linearity)
\item
  \textbf{y} = output
\end{itemize}

A single perceptron is a linear classifier --- it can only separate linearly separable classes. Famous failure: XOR problem (fixed by adding hidden layers).

\textbf{Interview takeaway:} A perceptron without a non-linear activation is just linear regression. Non-linearity is what makes neural nets powerful.

\subsection{Multilayer Perceptron (MLP)}\label{multilayer-perceptron-mlp}

Stack multiple perceptrons in layers:

\setcodelang{}

\begin{verbatim}
Input Layer → Hidden Layer 1 → Hidden Layer 2 → ... → Output Layer
\end{verbatim}

Each layer: \texttt{h\ =\ f(W·h\_prev\ +\ b)}

\textbf{Why layers matter:} Each layer learns a progressively more abstract representation:

\begin{itemize}
\tightlist
\item
  Layer 1 might detect low-level features (edges, n-gram frequencies)
\item
  Layer 2 combines those into mid-level patterns
\item
  Layer N learns task-specific abstractions
\end{itemize}

\textbf{Parameters:} \texttt{sum\ over\ layers\ of\ (n\_in\ *\ n\_out\ +\ n\_out)} --- weights + biases per layer.

\subsection{Activation Functions}\label{activation-functions}

\visualfigure{../figures/aiml-03-dl/01-activations.pdf}{Nonlinearities let networks model non-linear functions. ReLU/GELU dominate deep nets (cheap, no saturation for x>0); sigmoid/tanh saturate and can vanish gradients.}

Activation functions introduce \textbf{non-linearity} --- without them, stacking layers collapses to a single linear transformation (\texttt{W3\ *\ W2\ *\ W1\ *\ x\ =\ W\_combined\ *\ x}).

\subsubsection{Sigmoid}\label{sigmoid}

\setcodelang{}

\begin{verbatim}
σ(x) = 1 / (1 + e^(-x))    range: (0, 1)
\end{verbatim}

\begin{itemize}
\tightlist
\item
  Outputs are probability-like (range 0--1)
\item
  \textbf{Problem:} Saturates at extremes \(\rightarrow\) gradients \(\approx\) 0 \(\rightarrow\) \textbf{vanishing gradient}
\item
  \textbf{Problem:} Outputs not zero-centered \(\rightarrow\) zig-zag gradient updates
\item
  \textbf{Use:} Binary output layer only (never in hidden layers today)
\end{itemize}

\subsubsection{Tanh}\label{tanh}

\setcodelang{}

\begin{verbatim}
tanh(x) = (e^x - e^(-x)) / (e^x + e^(-x))    range: (-1, 1)
\end{verbatim}

\begin{itemize}
\tightlist
\item
  Zero-centered (better than sigmoid for hidden layers)
\item
  Still saturates \(\rightarrow\) vanishing gradients
\item
  \textbf{Use:} RNN hidden states (historically common)
\end{itemize}

\subsubsection{ReLU (Rectified Linear Unit)}\label{relu-rectified-linear-unit}

\setcodelang{}

\begin{verbatim}
ReLU(x) = max(0, x)
\end{verbatim}

\begin{itemize}
\tightlist
\item
  \textbf{No saturation for x \textgreater{} 0} \(\rightarrow\) gradients flow freely
\item
  Computationally cheap (just a threshold)
\item
  \textbf{Problem: Dead neurons} --- if a neuron\textquotesingle s input is always negative, it never activates, gradient = 0 forever, it "dies"
\item
  \textbf{Use:} Default for hidden layers in most modern networks
\end{itemize}

\subsubsection{Leaky ReLU}\label{leaky-relu}

\setcodelang{}

\begin{verbatim}
LeakyReLU(x) = max(0.01*x, x)
\end{verbatim}

\begin{itemize}
\tightlist
\item
  Small negative slope (0.01) prevents dead neurons
\item
  Small computational overhead vs ReLU
\item
  \textbf{Use:} When dead neuron problem is observed
\end{itemize}

\subsubsection{GELU (Gaussian Error Linear Unit)}\label{gelu-gaussian-error-linear-unit}

\setcodelang{}

\begin{verbatim}
GELU(x) ≈ 0.5 * x * (1 + tanh(√(2/π) * (x + 0.044715 * x³)))
\end{verbatim}

\begin{itemize}
\tightlist
\item
  Smooth, non-monotonic around zero
\item
  Better gradient properties than ReLU in practice
\item
  \textbf{Use:} Transformers (BERT, GPT), modern architectures
\item
  \textbf{Intuition:} Stochastically gates input based on its magnitude
\end{itemize}

\subsubsection{Softmax}\label{softmax}

\setcodelang{}

\begin{verbatim}
softmax(x_i) = e^(x_i) / sum_j(e^(x_j))
\end{verbatim}

\begin{itemize}
\tightlist
\item
  Outputs sum to 1 \(\rightarrow\) probability distribution over classes
\item
  Used only in \textbf{output layer} for multi-class classification
\item
  Pairs with cross-entropy loss
\item
  \textbf{Numerical stability trick:} subtract max(x) before exponentiating
\end{itemize}

\subsubsection{Activation Function Comparison Table}\label{activation-function-comparison-table}

\begin{longtable}[]{@{}
  >{\raggedright\arraybackslash}p{(\columnwidth - 10\tabcolsep) * \real{0.1181}}
  >{\raggedright\arraybackslash}p{(\columnwidth - 10\tabcolsep) * \real{0.1644}}
  >{\raggedright\arraybackslash}p{(\columnwidth - 10\tabcolsep) * \real{0.1536}}
  >{\raggedright\arraybackslash}p{(\columnwidth - 10\tabcolsep) * \real{0.1654}}
  >{\raggedright\arraybackslash}p{(\columnwidth - 10\tabcolsep) * \real{0.1418}}
  >{\raggedright\arraybackslash}p{(\columnwidth - 10\tabcolsep) * \real{0.2568}}@{}}
\toprule\noalign{}
\rowcolor{acc!16}
\begin{minipage}[b]{\linewidth}\raggedright
Activation
\end{minipage} & \begin{minipage}[b]{\linewidth}\raggedright
Range
\end{minipage} & \begin{minipage}[b]{\linewidth}\raggedright
Zero-Centered
\end{minipage} & \begin{minipage}[b]{\linewidth}\raggedright
Vanishing Grad
\end{minipage} & \begin{minipage}[b]{\linewidth}\raggedright
Dead Neurons
\end{minipage} & \begin{minipage}[b]{\linewidth}\raggedright
Primary Use
\end{minipage} \\
\midrule\noalign{}
\endhead
\bottomrule\noalign{}
\endlastfoot
Sigmoid & (0,1) & No & Yes (severe) & No & Binary output \\
\rowcolor{zebra}
Tanh & (-1,1) & Yes & Yes (mild) & No & RNN hidden states \\
ReLU & {[}0,\(\infty\)) & No & No (pos side) & Yes & Hidden layers (default) \\
\rowcolor{zebra}
Leaky ReLU & (-\(\infty\),\(\infty\)) & No & No & No & When dead neurons occur \\
GELU & (-\(\infty\),\(\infty\)) & No & No & No & Transformers, modern nets \\
\rowcolor{zebra}
Softmax & (0,1), sum=1 & No & N/A & N/A & Multiclass output \\
\end{longtable}

\textbf{Why non-linearity matters --- the collapse argument:}

\setcodelang{Python}

\begin{Shaded}
\begin{Highlighting}[]
\CommentTok{\# WITHOUT activation (just linear layers):}
\NormalTok{h1 }\OperatorTok{=}\NormalTok{ W1 }\OperatorTok{@}\NormalTok{ x }\OperatorTok{+}\NormalTok{ b1}
\NormalTok{h2 }\OperatorTok{=}\NormalTok{ W2 }\OperatorTok{@}\NormalTok{ h1 }\OperatorTok{+}\NormalTok{ b2}
\CommentTok{\# h2 = W2 @ (W1 @ x + b1) + b2}
\CommentTok{\#    = (W2@W1) @ x + (W2@b1 + b2)}
\CommentTok{\#    = W\_eff @ x + b\_eff   \textless{}{-} still linear!}

\CommentTok{\# WITH activation:}
\NormalTok{h1 }\OperatorTok{=}\NormalTok{ relu(W1 }\OperatorTok{@}\NormalTok{ x }\OperatorTok{+}\NormalTok{ b1)}
\NormalTok{h2 }\OperatorTok{=}\NormalTok{ relu(W2 }\OperatorTok{@}\NormalTok{ h1 }\OperatorTok{+}\NormalTok{ b2)}
\CommentTok{\# Cannot be collapsed {-}\textgreater{} genuinely non{-}linear function}
\end{Highlighting}
\end{Shaded}

\subsection{Forward Propagation}\label{forward-propagation}

The \textbf{forward pass} computes the network\textquotesingle s output from input to prediction:

\setcodelang{Python}

\begin{Shaded}
\begin{Highlighting}[]
\KeywordTok{def}\NormalTok{ forward(x, layers):}
    \CommentTok{"""}
\CommentTok{    x: input tensor [batch\_size, input\_dim]}
\CommentTok{    layers: list of (W, b, activation) tuples}
\CommentTok{    """}
\NormalTok{    h }\OperatorTok{=}\NormalTok{ x}
    \ControlFlowTok{for}\NormalTok{ W, b, act }\KeywordTok{in}\NormalTok{ layers:}
\NormalTok{        z }\OperatorTok{=}\NormalTok{ h }\OperatorTok{@}\NormalTok{ W }\OperatorTok{+}\NormalTok{ b        }\CommentTok{\# linear transformation}
\NormalTok{        h }\OperatorTok{=}\NormalTok{ act(z)           }\CommentTok{\# non{-}linear activation}
    \ControlFlowTok{return}\NormalTok{ h                 }\CommentTok{\# final output (logits or probabilities)}
\end{Highlighting}
\end{Shaded}

\textbf{Key tensors:}

\begin{itemize}
\tightlist
\item
  \texttt{z\^{}(l)\ =\ W\^{}(l)\ *\ a\^{}(l-1)\ +\ b\^{}(l)} --- pre-activation (logit)
\item
  \texttt{a\^{}(l)\ =\ f(z\^{}(l))} --- post-activation (activation)
\end{itemize}

Store both \texttt{z} and \texttt{a} for each layer --- needed in backprop.

\subsection{Loss Functions}\label{loss-functions}

The loss function measures how wrong the network\textquotesingle s predictions are. It must be \textbf{differentiable} so gradients can flow.

\subsubsection{Mean Squared Error (MSE) --- Regression}\label{mean-squared-error-mse--regression}

\setcodelang{Python}

\begin{Shaded}
\begin{Highlighting}[]
\NormalTok{MSE }\OperatorTok{=}\NormalTok{ (}\DecValTok{1}\OperatorTok{/}\NormalTok{n) }\OperatorTok{*}\NormalTok{ sum\_i( (y\_i }\OperatorTok{{-}}\NormalTok{ y\_hat\_i)}\OperatorTok{\^{}}\DecValTok{2}\NormalTok{ )}
\end{Highlighting}
\end{Shaded}

\begin{itemize}
\tightlist
\item
  Penalizes large errors quadratically
\item
  Sensitive to outliers (use MAE if outliers are common)
\item
  Gradient: \texttt{dL/dy\_hat\ =\ -2/n\ *\ (y\ -\ y\_hat)}
\end{itemize}

\subsubsection{Binary Cross-Entropy --- Binary Classification}\label{binary-cross-entropy--binary-classification}

\setcodelang{}

\begin{verbatim}
BCE = -( y * log(y_hat) + (1-y) * log(1-y_hat) )
\end{verbatim}

\begin{itemize}
\tightlist
\item
  Assumes \texttt{y\_hat\ =\ sigmoid(logit)} \(\in\) (0,1)
\item
  Penalizes confident wrong predictions heavily (log of near-zero \(\rightarrow\) -\(\infty\))
\end{itemize}

\subsubsection{Categorical Cross-Entropy --- Multiclass}\label{categorical-cross-entropy--multiclass}

\setcodelang{}

\begin{verbatim}
CE = -sum_c( y_c * log(y_hat_c) )
\end{verbatim}

\begin{itemize}
\tightlist
\item
  \texttt{y} is one-hot encoded, \texttt{y\_hat\ =\ softmax(logits)}
\item
  In practice: only the log of the predicted probability for the true class matters
\item
  Combined with softmax \(\rightarrow\) "softmax + cross-entropy" (numerically stable together)
\end{itemize}

\subsubsection{Cross-Entropy vs MSE for Classification}\label{cross-entropy-vs-mse-for-classification}

\begin{longtable}[]{@{}
  >{\raggedright\arraybackslash}p{(\columnwidth - 4\tabcolsep) * \real{0.2650}}
  >{\raggedright\arraybackslash}p{(\columnwidth - 4\tabcolsep) * \real{0.3505}}
  >{\raggedright\arraybackslash}p{(\columnwidth - 4\tabcolsep) * \real{0.3846}}@{}}
\toprule\noalign{}
\rowcolor{acc!16}
\begin{minipage}[b]{\linewidth}\raggedright
\end{minipage} & \begin{minipage}[b]{\linewidth}\raggedright
Cross-Entropy
\end{minipage} & \begin{minipage}[b]{\linewidth}\raggedright
MSE
\end{minipage} \\
\midrule\noalign{}
\endhead
\bottomrule\noalign{}
\endlastfoot
Appropriate for & Classification & Regression \\
\rowcolor{zebra}
Gradient behavior & Strong gradient even when very wrong & Weak gradient when output is saturated \\
Output activation & Softmax/Sigmoid & Linear \\
\rowcolor{zebra}
Why CE is better for classification & Log loss aligns with MLE for categorical distributions & Gradients vanish with sigmoid/softmax output when error is large \\
\end{longtable}

\subsection{Backpropagation \& the Chain Rule}\label{backpropagation--the-chain-rule}

Backpropagation is \textbf{automatic differentiation applied to a computational graph} --- computing \texttt{dL/dW} for every weight \texttt{W} using the chain rule of calculus.

\textbf{The chain rule:}

\setcodelang{Python}

\begin{Shaded}
\begin{Highlighting}[]
\NormalTok{If y }\OperatorTok{=}\NormalTok{ f(g(x)), then dy}\OperatorTok{/}\NormalTok{dx }\OperatorTok{=}\NormalTok{ dy}\OperatorTok{/}\NormalTok{dg }\OperatorTok{*}\NormalTok{ dg}\OperatorTok{/}\NormalTok{dx}

\NormalTok{More generally }\ControlFlowTok{for}\NormalTok{ composed functions:}
\NormalTok{dL}\OperatorTok{/}\NormalTok{dW}\OperatorTok{\^{}}\NormalTok{(l) }\OperatorTok{=}\NormalTok{ dL}\OperatorTok{/}\NormalTok{da}\OperatorTok{\^{}}\NormalTok{(l) }\OperatorTok{*}\NormalTok{ da}\OperatorTok{\^{}}\NormalTok{(l)}\OperatorTok{/}\NormalTok{dz}\OperatorTok{\^{}}\NormalTok{(l) }\OperatorTok{*}\NormalTok{ dz}\OperatorTok{\^{}}\NormalTok{(l)}\OperatorTok{/}\NormalTok{dW}\OperatorTok{\^{}}\NormalTok{(l)}
\end{Highlighting}
\end{Shaded}

\textbf{Backprop algorithm:}

\setcodelang{Python}

\begin{Shaded}
\begin{Highlighting}[]
\KeywordTok{def}\NormalTok{ backward(loss, layers, stored\_activations):}
    \CommentTok{"""}
\CommentTok{    Compute gradients for all weights using chain rule.}
\CommentTok{    stored\_activations: (z\^{}l, a\^{}l) for each layer from forward pass}
\CommentTok{    """}
\NormalTok{    dL\_da }\OperatorTok{=}\NormalTok{ d\_loss\_d\_output(loss)           }\CommentTok{\# gradient of loss w.r.t. output}
    
    \ControlFlowTok{for}\NormalTok{ l }\KeywordTok{in} \BuiltInTok{reversed}\NormalTok{(}\BuiltInTok{range}\NormalTok{(}\BuiltInTok{len}\NormalTok{(layers))):}
\NormalTok{        z\_l, a\_prev }\OperatorTok{=}\NormalTok{ stored\_activations[l]}
\NormalTok{        W\_l, b\_l, act\_l }\OperatorTok{=}\NormalTok{ layers[l]}
        
        \CommentTok{\# Step 1: gradient through activation function}
\NormalTok{        dL\_dz }\OperatorTok{=}\NormalTok{ dL\_da }\OperatorTok{*}\NormalTok{ act\_l.derivative(z\_l)   }\CommentTok{\# element{-}wise}
        
        \CommentTok{\# Step 2: gradient w.r.t. weights and bias}
\NormalTok{        dL\_dW }\OperatorTok{=}\NormalTok{ a\_prev.T }\OperatorTok{@}\NormalTok{ dL\_dz                }\CommentTok{\# [n\_in, batch] @ [batch, n\_out]}
\NormalTok{        dL\_db }\OperatorTok{=}\NormalTok{ dL\_dz.}\BuiltInTok{sum}\NormalTok{(axis}\OperatorTok{=}\DecValTok{0}\NormalTok{)               }\CommentTok{\# sum over batch}
        
        \CommentTok{\# Step 3: gradient to pass to previous layer}
\NormalTok{        dL\_da }\OperatorTok{=}\NormalTok{ dL\_dz }\OperatorTok{@}\NormalTok{ W\_l.T                   }\CommentTok{\# [batch, n\_in]}
        
        \CommentTok{\# Store for optimizer}
\NormalTok{        gradients[l] }\OperatorTok{=}\NormalTok{ (dL\_dW, dL\_db)}
    
    \ControlFlowTok{return}\NormalTok{ gradients}
\end{Highlighting}
\end{Shaded}

\textbf{Key insight:} The gradient "flows backward" through each layer. Each layer contributes a factor to the chain --- if any factor is near zero (saturated activation), the gradient vanishes.

\textbf{Interview takeaway:} Backprop is just the chain rule applied layer by layer. The "backward" graph is the transpose of the forward graph.

\subsection{Optimizers}\label{optimizers}

Optimizers determine \textbf{how} gradient information is used to update weights.

\subsubsection{SGD (Stochastic Gradient Descent)}\label{sgd-stochastic-gradient-descent}

\visualfigure{../figures/aiml-03-dl/02-gradient-descent.pdf}{Each step moves opposite the gradient. Elongated (ill-conditioned) valleys make vanilla gradient descent zig-zag across the steep axis — motivating momentum, RMSProp and Adam.}

\setcodelang{Python}

\begin{Shaded}
\begin{Highlighting}[]
\NormalTok{W }\OperatorTok{=}\NormalTok{ W }\OperatorTok{{-}}\NormalTok{ lr }\OperatorTok{*}\NormalTok{ dL}\OperatorTok{/}\NormalTok{dW}
\end{Highlighting}
\end{Shaded}

\begin{itemize}
\tightlist
\item
  Simple, well-understood, good generalization
\item
  \textbf{Problem:} Single fixed learning rate, slow convergence, oscillates in narrow ravines
\item
  \textbf{Stochastic} = uses mini-batches (random subset), not full dataset --- noisy but faster
\end{itemize}

\subsubsection{SGD with Momentum}\label{sgd-with-momentum}

\setcodelang{Python}

\begin{Shaded}
\begin{Highlighting}[]
\NormalTok{v\_t }\OperatorTok{=}\NormalTok{ β }\OperatorTok{*}\NormalTok{ v\_\{t}\OperatorTok{{-}}\DecValTok{1}\NormalTok{\} }\OperatorTok{+}\NormalTok{ (}\DecValTok{1}\OperatorTok{{-}}\NormalTok{β) }\OperatorTok{*}\NormalTok{ dL}\OperatorTok{/}\NormalTok{dW\_t}
\NormalTok{W }\OperatorTok{=}\NormalTok{ W }\OperatorTok{{-}}\NormalTok{ lr }\OperatorTok{*}\NormalTok{ v\_t}
\end{Highlighting}
\end{Shaded}

\begin{itemize}
\tightlist
\item
  \texttt{v\_t} = velocity (exponential moving average of gradients)
\item
  \texttt{β} typically 0.9
\item
  \textbf{Effect:} Accelerates in consistent gradient directions, dampens oscillations
\item
  \textbf{Analogy:} A ball rolling downhill --- it builds momentum in the right direction
\end{itemize}

\subsubsection{RMSprop}\label{rmsprop}

\setcodelang{Python}

\begin{Shaded}
\begin{Highlighting}[]
\NormalTok{s\_t }\OperatorTok{=}\NormalTok{ β }\OperatorTok{*}\NormalTok{ s\_\{t}\OperatorTok{{-}}\DecValTok{1}\NormalTok{\} }\OperatorTok{+}\NormalTok{ (}\DecValTok{1}\OperatorTok{{-}}\NormalTok{β) }\OperatorTok{*}\NormalTok{ (dL}\OperatorTok{/}\NormalTok{dW\_t)²}
\NormalTok{W }\OperatorTok{=}\NormalTok{ W }\OperatorTok{{-}}\NormalTok{ (lr }\OperatorTok{/}\NormalTok{ sqrt(s\_t }\OperatorTok{+}\NormalTok{ ε)) }\OperatorTok{*}\NormalTok{ dL}\OperatorTok{/}\NormalTok{dW\_t}
\end{Highlighting}
\end{Shaded}

\begin{itemize}
\tightlist
\item
  Adapts learning rate \textbf{per parameter} based on recent gradient magnitudes
\item
  Parameters with large gradients get smaller effective LR
\item
  Good for RNNs and non-stationary objectives
\end{itemize}

\subsubsection{Adam (Adaptive Moment Estimation)}\label{adam-adaptive-moment-estimation}

\setcodelang{Python}

\begin{Shaded}
\begin{Highlighting}[]
\NormalTok{m\_t }\OperatorTok{=}\NormalTok{ β}\DecValTok{1} \OperatorTok{*}\NormalTok{ m\_\{t}\OperatorTok{{-}}\DecValTok{1}\NormalTok{\} }\OperatorTok{+}\NormalTok{ (}\DecValTok{1}\OperatorTok{{-}}\NormalTok{β}\DecValTok{1}\NormalTok{) }\OperatorTok{*}\NormalTok{ g\_t          }\CommentTok{\# 1st moment (mean)}
\NormalTok{v\_t }\OperatorTok{=}\NormalTok{ β}\DecValTok{2} \OperatorTok{*}\NormalTok{ v\_\{t}\OperatorTok{{-}}\DecValTok{1}\NormalTok{\} }\OperatorTok{+}\NormalTok{ (}\DecValTok{1}\OperatorTok{{-}}\NormalTok{β}\DecValTok{2}\NormalTok{) }\OperatorTok{*}\NormalTok{ g\_t²         }\CommentTok{\# 2nd moment (uncentered variance)}

\NormalTok{m\_hat }\OperatorTok{=}\NormalTok{ m\_t }\OperatorTok{/}\NormalTok{ (}\DecValTok{1} \OperatorTok{{-}}\NormalTok{ β}\DecValTok{1}\OperatorTok{\^{}}\NormalTok{t)                    }\CommentTok{\# bias correction}
\NormalTok{v\_hat }\OperatorTok{=}\NormalTok{ v\_t }\OperatorTok{/}\NormalTok{ (}\DecValTok{1} \OperatorTok{{-}}\NormalTok{ β}\DecValTok{2}\OperatorTok{\^{}}\NormalTok{t)                    }\CommentTok{\# bias correction}

\NormalTok{W }\OperatorTok{=}\NormalTok{ W }\OperatorTok{{-}}\NormalTok{ lr }\OperatorTok{*}\NormalTok{ m\_hat }\OperatorTok{/}\NormalTok{ (sqrt(v\_hat) }\OperatorTok{+}\NormalTok{ ε)}
\end{Highlighting}
\end{Shaded}

\begin{itemize}
\tightlist
\item
  \textbf{\(\beta\)1=0.9, \(\beta\)2=0.999, \(\varepsilon\)=1e-8} (defaults)
\item
  Combines momentum (1st moment) + per-parameter LR adaptation (2nd moment)
\item
  Bias correction handles initialization bias (m\_0 = v\_0 = 0)
\item
  \textbf{Default choice for most tasks}
\item
  \textbf{Weakness:} Can generalize slightly worse than SGD+momentum on some tasks (Adam converges faster but to a slightly different minimum)
\end{itemize}

\subsubsection{Learning Rate Schedules}\label{learning-rate-schedules}

\begin{longtable}[]{@{}
  >{\raggedright\arraybackslash}p{(\columnwidth - 4\tabcolsep) * \real{0.1799}}
  >{\raggedright\arraybackslash}p{(\columnwidth - 4\tabcolsep) * \real{0.4709}}
  >{\raggedright\arraybackslash}p{(\columnwidth - 4\tabcolsep) * \real{0.3492}}@{}}
\toprule\noalign{}
\rowcolor{acc!16}
\begin{minipage}[b]{\linewidth}\raggedright
Schedule
\end{minipage} & \begin{minipage}[b]{\linewidth}\raggedright
Formula / Description
\end{minipage} & \begin{minipage}[b]{\linewidth}\raggedright
Use Case
\end{minipage} \\
\midrule\noalign{}
\endhead
\bottomrule\noalign{}
\endlastfoot
Step Decay & LR *= \(\gamma\) every N epochs & Simple, works well \\
\rowcolor{zebra}
Cosine Annealing & LR = LR\_min + 0.5*(LR\_max-LR\_min)*(1+cos(\(\pi\)*t/T)) & Transformers, fine-tuning \\
Warmup + Decay & Linear warmup for K steps, then decay & Large model training (GPT, BERT) \\
\rowcolor{zebra}
Cyclic LR & Oscillate between LR\_min and LR\_max & Escape local minima \\
ReduceLROnPlateau & Reduce LR when metric stops improving & When you don\textquotesingle t know when to decay \\
\end{longtable}

\textbf{Why warmup?} Large models are unstable early in training --- gradients are noisy, momentum is uninitialized. Starting with a small LR and ramping up prevents early divergence.

\subsubsection{Optimizer Comparison}\label{optimizer-comparison}

\begin{longtable}[]{@{}
  >{\raggedright\arraybackslash}p{(\columnwidth - 8\tabcolsep) * \real{0.1503}}
  >{\raggedright\arraybackslash}p{(\columnwidth - 8\tabcolsep) * \real{0.1440}}
  >{\raggedright\arraybackslash}p{(\columnwidth - 8\tabcolsep) * \real{0.1296}}
  >{\raggedright\arraybackslash}p{(\columnwidth - 8\tabcolsep) * \real{0.1127}}
  >{\raggedright\arraybackslash}p{(\columnwidth - 8\tabcolsep) * \real{0.4634}}@{}}
\toprule\noalign{}
\rowcolor{acc!16}
\begin{minipage}[b]{\linewidth}\raggedright
Optimizer
\end{minipage} & \begin{minipage}[b]{\linewidth}\raggedright
Adapts LR?
\end{minipage} & \begin{minipage}[b]{\linewidth}\raggedright
Momentum?
\end{minipage} & \begin{minipage}[b]{\linewidth}\raggedright
Memory
\end{minipage} & \begin{minipage}[b]{\linewidth}\raggedright
Best For
\end{minipage} \\
\midrule\noalign{}
\endhead
\bottomrule\noalign{}
\endlastfoot
SGD & No & No & 1x params & Final fine-tuning, convex problems \\
\rowcolor{zebra}
SGD+Momentum & No & Yes & 2x params & Vision models (ResNet, etc.) \\
RMSprop & Yes & No & 2x params & RNNs, non-stationary \\
\rowcolor{zebra}
Adam & Yes & Yes & 3x params & Default: most deep learning tasks \\
AdamW & Yes & Yes & 3x params & Transformers (decouples weight decay) \\
\end{longtable}

\subsection{Vanishing \& Exploding Gradients}\label{vanishing--exploding-gradients}

\textbf{Vanishing gradients:} During backprop, gradients become exponentially small as they propagate through many layers \(\rightarrow\) early layers learn very slowly or not at all.

\textbf{Root cause:} Chain rule multiplies many numbers together. If each factor \textless{} 1 (e.g., sigmoid derivative max = 0.25), the product shrinks exponentially with depth.

\setcodelang{}

\begin{verbatim}
dL/dW^(1) = dL/da^(L) * (da^(L)/dz^(L)) * ... * (da^(2)/dz^(2)) * (da^(1)/dz^(1))
           = product of L terms, each ≤ 0.25 for sigmoid
           ≈ 0.25^L → 0 as L → ∞
\end{verbatim}

\textbf{Exploding gradients:} Same mechanism but factors \textgreater{} 1 \(\rightarrow\) gradients blow up \(\rightarrow\) NaN weights.

\textbf{Solutions:}

\begin{longtable}[]{@{}
  >{\raggedright\arraybackslash}p{(\columnwidth - 4\tabcolsep) * \real{0.1066}}
  >{\raggedright\arraybackslash}p{(\columnwidth - 4\tabcolsep) * \real{0.3555}}
  >{\raggedright\arraybackslash}p{(\columnwidth - 4\tabcolsep) * \real{0.5379}}@{}}
\toprule\noalign{}
\rowcolor{acc!16}
\begin{minipage}[b]{\linewidth}\raggedright
Problem
\end{minipage} & \begin{minipage}[b]{\linewidth}\raggedright
Solution
\end{minipage} & \begin{minipage}[b]{\linewidth}\raggedright
Mechanism
\end{minipage} \\
\midrule\noalign{}
\endhead
\bottomrule\noalign{}
\endlastfoot
Vanishing & ReLU activation & Derivative = 1 for x \textgreater{} 0, doesn\textquotesingle t shrink \\
\rowcolor{zebra}
Vanishing & Residual connections (ResNets) & Gradient highway: \texttt{dL/dW\^{}(l)} gets additive term of 1 \\
Vanishing & Better initialization & Xavier/He --- keep gradient magnitude stable \\
\rowcolor{zebra}
Vanishing & Batch normalization & Normalizes activations, keeps them in active range \\
Exploding & Gradient clipping & `g = g * (max\_norm / \\
\rowcolor{zebra}
Both & LSTM/GRU gating & Additive cell state update instead of multiplicative \\
\end{longtable}

\textbf{Interview takeaway:} ReLU + residual connections + BatchNorm together largely solved vanishing gradients for feedforward nets. LSTMs solved it for RNNs.

\subsection{Weight Initialization}\label{weight-initialization}

Poor initialization \(\rightarrow\) dead neurons, vanishing/exploding gradients from step 1 (before training even starts).

\textbf{Goal:} Keep the variance of activations and gradients approximately constant across layers.

\subsubsection{Xavier/Glorot Initialization (for Tanh/Sigmoid)}\label{xavierglorot-initialization-for-tanhsigmoid}

\setcodelang{}

\begin{verbatim}
W ~ Uniform(-√(6/(n_in + n_out)), √(6/(n_in + n_out)))
# or equivalently:
W ~ Normal(0, √(2/(n_in + n_out)))
\end{verbatim}

\begin{itemize}
\tightlist
\item
  Derived by requiring \texttt{Var(output)\ =\ Var(input)} for linear activations
\item
  Works well with tanh and sigmoid
\end{itemize}

\subsubsection{He Initialization (for ReLU)}\label{he-initialization-for-relu}

\setcodelang{}

\begin{verbatim}
W ~ Normal(0, √(2/n_in))
\end{verbatim}

\begin{itemize}
\tightlist
\item
  Accounts for ReLU zeroing out half the inputs (halves the effective variance)
\item
  \textbf{Use He init with ReLU networks}
\end{itemize}

\subsubsection{Why initialization matters}\label{why-initialization-matters}

\setcodelang{}

\begin{verbatim}
# Bad init: all zeros
W = 0 → all neurons compute same output → same gradient → stay identical forever
# "Symmetry breaking" is essential — init must be random

# Bad init: too large
activations → saturate → vanishing gradients from step 1

# Bad init: too small  
activations → near zero → gradients near zero → barely learns
\end{verbatim}

\textbf{Interview takeaway:} Xavier for tanh/sigmoid, He for ReLU. Always random (never all-zeros for weights --- zeros for biases is fine).

\subsection{Batch Normalization \& Layer Normalization}\label{batch-normalization--layer-normalization}

\subsubsection{Batch Normalization (BatchNorm)}\label{batch-normalization-batchnorm}

Normalize activations \textbf{across the batch dimension} for each feature:

\setcodelang{Python}

\begin{Shaded}
\begin{Highlighting}[]
\NormalTok{μ\_B }\OperatorTok{=}\NormalTok{ (}\DecValTok{1}\OperatorTok{/}\NormalTok{m) }\OperatorTok{*}\NormalTok{ sum\_i(x\_i)          }\CommentTok{\# batch mean}
\NormalTok{σ²\_B }\OperatorTok{=}\NormalTok{ (}\DecValTok{1}\OperatorTok{/}\NormalTok{m) }\OperatorTok{*}\NormalTok{ sum\_i((x\_i }\OperatorTok{{-}}\NormalTok{ μ\_B)²) }\CommentTok{\# batch variance}

\NormalTok{x\_hat\_i }\OperatorTok{=}\NormalTok{ (x\_i }\OperatorTok{{-}}\NormalTok{ μ\_B) }\OperatorTok{/}\NormalTok{ sqrt(σ²\_B }\OperatorTok{+}\NormalTok{ ε)   }\CommentTok{\# normalize}

\NormalTok{y\_i }\OperatorTok{=}\NormalTok{ γ }\OperatorTok{*}\NormalTok{ x\_hat\_i }\OperatorTok{+}\NormalTok{ β             }\CommentTok{\# learnable scale and shift}
\end{Highlighting}
\end{Shaded}

\begin{itemize}
\tightlist
\item
  \texttt{γ,\ β} are learned parameters (lets network un-normalize if needed)
\item
  Applied \textbf{before} the activation function (or after --- debate exists)
\item
  At inference: use running mean/variance tracked during training (not batch stats)
\end{itemize}

\textbf{Benefits:}

\begin{enumerate}
\def\labelenumi{\arabic{enumi}.}
\tightlist
\item
  Reduces internal covariate shift
\item
  Allows higher learning rates
\item
  Acts as slight regularizer (noise from batch statistics)
\item
  Reduces sensitivity to initialization
\end{enumerate}

\textbf{Problems with BatchNorm:}

\begin{itemize}
\tightlist
\item
  Batch size dependent --- fails for batch\_size=1
\item
  Can\textquotesingle t be used as-is for RNNs (variable-length sequences)
\item
  Creates dependency between samples in a batch
\end{itemize}

\subsubsection{Layer Normalization (LayerNorm)}\label{layer-normalization-layernorm}

Normalize \textbf{across the feature dimension} for each sample:

\setcodelang{Python}

\begin{Shaded}
\begin{Highlighting}[]
\NormalTok{μ }\OperatorTok{=}\NormalTok{ (}\DecValTok{1}\OperatorTok{/}\NormalTok{H) }\OperatorTok{*}\NormalTok{ sum\_j(x\_j)            }\CommentTok{\# mean over features}
\NormalTok{σ² }\OperatorTok{=}\NormalTok{ (}\DecValTok{1}\OperatorTok{/}\NormalTok{H) }\OperatorTok{*}\NormalTok{ sum\_j((x\_j }\OperatorTok{{-}}\NormalTok{ μ)²)   }\CommentTok{\# variance over features}

\NormalTok{x\_hat }\OperatorTok{=}\NormalTok{ (x }\OperatorTok{{-}}\NormalTok{ μ) }\OperatorTok{/}\NormalTok{ sqrt(σ² }\OperatorTok{+}\NormalTok{ ε)}
\NormalTok{y }\OperatorTok{=}\NormalTok{ γ }\OperatorTok{*}\NormalTok{ x\_hat }\OperatorTok{+}\NormalTok{ β}
\end{Highlighting}
\end{Shaded}

\begin{itemize}
\tightlist
\item
  Independent of batch size
\item
  Same behavior at train and inference time
\item
  \textbf{Works for RNNs and Transformers}
\item
  Default in all Transformer architectures (BERT, GPT, etc.)
\end{itemize}

\subsubsection{BatchNorm vs LayerNorm}\label{batchnorm-vs-layernorm}

\begin{longtable}[]{@{}
  >{\raggedright\arraybackslash}p{(\columnwidth - 4\tabcolsep) * \real{0.3418}}
  >{\raggedright\arraybackslash}p{(\columnwidth - 4\tabcolsep) * \real{0.3165}}
  >{\raggedright\arraybackslash}p{(\columnwidth - 4\tabcolsep) * \real{0.3418}}@{}}
\toprule\noalign{}
\rowcolor{acc!16}
\begin{minipage}[b]{\linewidth}\raggedright
\end{minipage} & \begin{minipage}[b]{\linewidth}\raggedright
BatchNorm
\end{minipage} & \begin{minipage}[b]{\linewidth}\raggedright
LayerNorm
\end{minipage} \\
\midrule\noalign{}
\endhead
\bottomrule\noalign{}
\endlastfoot
Normalization axis & Across batch, per feature & Across features, per sample \\
\rowcolor{zebra}
Batch size dependent & Yes & No \\
Works for RNNs/Transformers & No & Yes \\
\rowcolor{zebra}
Used in & CNNs, MLPs & RNNs, Transformers \\
Train/inference difference & Yes (running stats) & No \\
\rowcolor{zebra}
Acts as regularizer & Yes (batch noise) & Minimal \\
\end{longtable}

\subsection{Dropout \& Regularization in Deep Learning}\label{dropout--regularization-in-deep-learning}

\subsubsection{Dropout}\label{dropout}

During training, randomly set each neuron\textquotesingle s output to \textbf{zero with probability p} (typically p=0.2 to 0.5):

\setcodelang{Python}

\begin{Shaded}
\begin{Highlighting}[]
\KeywordTok{def}\NormalTok{ dropout(x, p, training}\OperatorTok{=}\VariableTok{True}\NormalTok{):}
    \ControlFlowTok{if} \KeywordTok{not}\NormalTok{ training:}
        \ControlFlowTok{return}\NormalTok{ x                          }\CommentTok{\# No dropout at inference}
\NormalTok{    mask }\OperatorTok{=}\NormalTok{ (torch.rand\_like(x) }\OperatorTok{\textgreater{}}\NormalTok{ p)      }\CommentTok{\# Bernoulli mask}
    \ControlFlowTok{return}\NormalTok{ x }\OperatorTok{*}\NormalTok{ mask }\OperatorTok{/}\NormalTok{ (}\DecValTok{1} \OperatorTok{{-}}\NormalTok{ p)            }\CommentTok{\# Scale to maintain expected value}
\end{Highlighting}
\end{Shaded}

\textbf{Why it works:}

\begin{enumerate}
\def\labelenumi{\arabic{enumi}.}
\tightlist
\item
  Prevents co-adaptation --- neurons can\textquotesingle t rely on specific other neurons
\item
  Equivalent to training an ensemble of \texttt{2\^{}n} different network architectures
\item
  At inference, full network with scaled weights approximates ensemble average
\end{enumerate}

\textbf{Where to apply:} After fully-connected layers, NOT typically after convolutions (or use spatial dropout for CNNs).

\subsubsection{Other DL Regularization Techniques}\label{other-dl-regularization-techniques}

\begin{longtable}[]{@{}
  >{\raggedright\arraybackslash}p{(\columnwidth - 4\tabcolsep) * \real{0.1640}}
  >{\raggedright\arraybackslash}p{(\columnwidth - 4\tabcolsep) * \real{0.4597}}
  >{\raggedright\arraybackslash}p{(\columnwidth - 4\tabcolsep) * \real{0.3763}}@{}}
\toprule\noalign{}
\rowcolor{acc!16}
\begin{minipage}[b]{\linewidth}\raggedright
Technique
\end{minipage} & \begin{minipage}[b]{\linewidth}\raggedright
Mechanism
\end{minipage} & \begin{minipage}[b]{\linewidth}\raggedright
When to Use
\end{minipage} \\
\midrule\noalign{}
\endhead
\bottomrule\noalign{}
\endlastfoot
L2 (Weight Decay) & \texttt{loss\ +=\ λ\ *\ sum(W²)} \(\rightarrow\) pulls weights toward zero & Always, as AdamW \\
\rowcolor{zebra}
L1 & `loss += λ * sum( & W \\
Dropout & Random neuron deactivation & FC layers, Transformers \\
\rowcolor{zebra}
Data Augmentation & Random crops, flips, color jitter & Image tasks \\
Early Stopping & Stop when val loss stops improving & All tasks \\
\rowcolor{zebra}
Label Smoothing & \texttt{y\ =\ (1-ε)*y\_onehot\ +\ ε/K} & Classification, prevents overconfidence \\
Batch Norm & Noise from batch statistics & CNNs \\
\end{longtable}

\subsection{Convolutional Neural Networks (CNNs)}\label{convolutional-neural-networks-cnns}

CNNs exploit \textbf{spatial structure} --- nearby pixels are correlated. Instead of connecting every input pixel to every neuron (which would be computationally infeasible for images), use \textbf{local, shared filters}.

\subsubsection{Core Components}\label{core-components}

\textbf{Convolution Operation:}

\setcodelang{}

\begin{verbatim}
(f * I)[i,j] = sum_m sum_n f[m,n] * I[i+m, j+n]
\end{verbatim}

\begin{itemize}
\tightlist
\item
  \textbf{Filter/Kernel:} Small weight matrix (e.g., 3\(\times\)3, 5\(\times\)5) that slides over input
\item
  \textbf{Feature map:} Output of applying one filter to the full input
\item
  \textbf{Multiple filters} \(\rightarrow\) multiple feature maps \(\rightarrow\) detect different patterns
\item
  \textbf{Shared weights:} Same filter weights applied at every position \(\rightarrow\) translation invariance + far fewer parameters than FC
\end{itemize}

\textbf{Stride:} How many pixels the filter moves each step.

\begin{itemize}
\tightlist
\item
  Stride 1: output \(\approx\) same size as input
\item
  Stride 2: halves spatial dimensions
\end{itemize}

\textbf{Padding:}

\begin{itemize}
\tightlist
\item
  "Valid": no padding \(\rightarrow\) output smaller than input
\item
  "Same": pad with zeros \(\rightarrow\) output same size as input
\end{itemize}

\textbf{Output size formula:}

\setcodelang{Python}

\begin{Shaded}
\begin{Highlighting}[]
\NormalTok{output\_size }\OperatorTok{=}\NormalTok{ floor((input\_size }\OperatorTok{{-}}\NormalTok{ kernel\_size }\OperatorTok{+} \DecValTok{2}\OperatorTok{*}\NormalTok{padding) }\OperatorTok{/}\NormalTok{ stride) }\OperatorTok{+} \DecValTok{1}
\end{Highlighting}
\end{Shaded}

\textbf{Pooling:}

\begin{itemize}
\tightlist
\item
  Max pooling: takes maximum in each local region \(\rightarrow\) spatial downsampling + translation invariance
\item
  Average pooling: takes average
\item
  Global Average Pooling (GAP): pools entire feature map to single value per channel
\end{itemize}

\textbf{Receptive Field:} The area of the original input that influences one output neuron. Deeper layers have larger receptive fields.

\subsubsection{Classic CNN Architectures}\label{classic-cnn-architectures}

\begin{longtable}[]{@{}
  >{\raggedright\arraybackslash}p{(\columnwidth - 6\tabcolsep) * \real{0.2417}}
  >{\raggedright\arraybackslash}p{(\columnwidth - 6\tabcolsep) * \real{0.0600}}
  >{\raggedright\arraybackslash}p{(\columnwidth - 6\tabcolsep) * \real{0.6120}}
  >{\raggedright\arraybackslash}p{(\columnwidth - 6\tabcolsep) * \real{0.0878}}@{}}
\toprule\noalign{}
\rowcolor{acc!16}
\begin{minipage}[b]{\linewidth}\raggedright
Architecture
\end{minipage} & \begin{minipage}[b]{\linewidth}\raggedright
Year
\end{minipage} & \begin{minipage}[b]{\linewidth}\raggedright
Key Innovation
\end{minipage} & \begin{minipage}[b]{\linewidth}\raggedright
Params
\end{minipage} \\
\midrule\noalign{}
\endhead
\bottomrule\noalign{}
\endlastfoot
LeNet-5 & 1998 & Conv-Pool-FC design & \textasciitilde60K \\
\rowcolor{zebra}
AlexNet & 2012 & Deep CNN + ReLU + Dropout + GPU & \textasciitilde60M \\
VGG-16 & 2014 & Uniform 3\(\times\)3 convs, very deep & \textasciitilde138M \\
\rowcolor{zebra}
GoogLeNet/Inception & 2014 & Inception modules (multi-scale filters) & \textasciitilde6.8M \\
ResNet-50 & 2015 & Residual connections \(\rightarrow\) very deep (50-152 layers) & \textasciitilde25M \\
\rowcolor{zebra}
EfficientNet & 2019 & Compound scaling (depth, width, resolution) & 5-66M \\
\end{longtable}

\textbf{ResNet key innovation --- skip connections:}

\setcodelang{Python}

\begin{Shaded}
\begin{Highlighting}[]
\NormalTok{h }\OperatorTok{=}\NormalTok{ F(x, W) }\OperatorTok{+}\NormalTok{ x    }\CommentTok{\# Instead of just h = F(x, W)}
\end{Highlighting}
\end{Shaded}

The \texttt{+\ x} creates a "gradient highway" --- backprop can flow directly through the identity connection, bypassing potentially vanishing layers.

\subsubsection{Parameter Count for Convolution Layer}\label{parameter-count-for-convolution-layer}

\setcodelang{Python}

\begin{Shaded}
\begin{Highlighting}[]
\NormalTok{params }\OperatorTok{=}\NormalTok{ (kernel\_h }\OperatorTok{*}\NormalTok{ kernel\_w }\OperatorTok{*}\NormalTok{ in\_channels }\OperatorTok{+} \DecValTok{1}\NormalTok{) }\OperatorTok{*}\NormalTok{ out\_channels}
\CommentTok{\#         weights per filter                       number of filters}
\CommentTok{\#         +1 for bias}
\end{Highlighting}
\end{Shaded}

Compare to FC layer: \texttt{params\ =\ in\_size\ *\ out\_size\ +\ out\_size}. For images, conv is massively more parameter-efficient.

\subsection{Recurrent Neural Networks (RNNs)}\label{recurrent-neural-networks-rnns}

RNNs process \textbf{sequential data} by maintaining a hidden state that carries information across time steps:

\setcodelang{Python}

\begin{Shaded}
\begin{Highlighting}[]
\KeywordTok{def}\NormalTok{ rnn\_step(x\_t, h\_prev, W\_xh, W\_hh, b\_h, W\_hy, b\_y):}
\NormalTok{    h\_t }\OperatorTok{=}\NormalTok{ tanh(W\_xh }\OperatorTok{@}\NormalTok{ x\_t }\OperatorTok{+}\NormalTok{ W\_hh }\OperatorTok{@}\NormalTok{ h\_prev }\OperatorTok{+}\NormalTok{ b\_h)  }\CommentTok{\# hidden state update}
\NormalTok{    y\_t }\OperatorTok{=}\NormalTok{ W\_hy }\OperatorTok{@}\NormalTok{ h\_t }\OperatorTok{+}\NormalTok{ b\_y                           }\CommentTok{\# output (optional per step)}
    \ControlFlowTok{return}\NormalTok{ h\_t, y\_t}
\end{Highlighting}
\end{Shaded}

\textbf{The same weights (W\_xh, W\_hh) are shared across all time steps} --- this is analogous to weight sharing in CNNs.

\textbf{Problems with vanilla RNNs:}

\begin{enumerate}
\def\labelenumi{\arabic{enumi}.}
\tightlist
\item
  \textbf{Vanishing gradients through time:} \texttt{dL/dh\_0} = product of Jacobians across T steps \(\rightarrow\) exponentially small for long sequences
\item
  \textbf{Exploding gradients:} Same mechanism in reverse
\item
  \textbf{Short-term memory:} Effectively forgets information from \textgreater10-20 steps back
\end{enumerate}

\subsubsection{LSTM (Long Short-Term Memory)}\label{lstm-long-short-term-memory}

LSTMs solve the vanishing gradient problem with a \textbf{cell state} (long-term memory) and \textbf{gating mechanisms}:

\setcodelang{Python}

\begin{Shaded}
\begin{Highlighting}[]
\CommentTok{\# Gates (all computed same way: sigmoid of linear combination)}
\NormalTok{i\_t }\OperatorTok{=}\NormalTok{ sigmoid(W\_i }\OperatorTok{@}\NormalTok{ [h\_\{t}\OperatorTok{{-}}\DecValTok{1}\NormalTok{\}, x\_t] }\OperatorTok{+}\NormalTok{ b\_i)   }\CommentTok{\# input gate: what to write}
\NormalTok{f\_t }\OperatorTok{=}\NormalTok{ sigmoid(W\_f }\OperatorTok{@}\NormalTok{ [h\_\{t}\OperatorTok{{-}}\DecValTok{1}\NormalTok{\}, x\_t] }\OperatorTok{+}\NormalTok{ b\_f)   }\CommentTok{\# forget gate: what to erase}
\NormalTok{o\_t }\OperatorTok{=}\NormalTok{ sigmoid(W\_o }\OperatorTok{@}\NormalTok{ [h\_\{t}\OperatorTok{{-}}\DecValTok{1}\NormalTok{\}, x\_t] }\OperatorTok{+}\NormalTok{ b\_o)   }\CommentTok{\# output gate: what to read}

\CommentTok{\# Candidate cell state}
\NormalTok{c\_hat\_t }\OperatorTok{=}\NormalTok{ tanh(W\_c }\OperatorTok{@}\NormalTok{ [h\_\{t}\OperatorTok{{-}}\DecValTok{1}\NormalTok{\}, x\_t] }\OperatorTok{+}\NormalTok{ b\_c)}

\CommentTok{\# Cell state update (ADDITIVE — preserves gradient flow!)}
\NormalTok{c\_t }\OperatorTok{=}\NormalTok{ f\_t }\OperatorTok{*}\NormalTok{ c\_\{t}\OperatorTok{{-}}\DecValTok{1}\NormalTok{\} }\OperatorTok{+}\NormalTok{ i\_t }\OperatorTok{*}\NormalTok{ c\_hat\_t}

\CommentTok{\# Hidden state}
\NormalTok{h\_t }\OperatorTok{=}\NormalTok{ o\_t }\OperatorTok{*}\NormalTok{ tanh(c\_t)}
\end{Highlighting}
\end{Shaded}

\textbf{Why LSTMs work:}

\begin{itemize}
\tightlist
\item
  \textbf{Additive cell update:} \texttt{c\_t\ =\ f\_t\ *\ c\_\{t-1\}\ +\ i\_t\ *\ c\_hat\_t} --- gradient can flow through addition without multiplication, analogous to ResNet skip connections
\item
  \textbf{Forget gate \(\approx\) 1:} If forget gate is near 1, cell state is preserved unchanged, gradient flows unimpeded
\item
  Separate \texttt{h\_t} (working memory) and \texttt{c\_t} (long-term memory)
\end{itemize}

\subsubsection{GRU (Gated Recurrent Unit)}\label{gru-gated-recurrent-unit}

Simplified LSTM with two gates instead of three:

\setcodelang{Python}

\begin{Shaded}
\begin{Highlighting}[]
\NormalTok{z\_t }\OperatorTok{=}\NormalTok{ sigmoid(W\_z }\OperatorTok{@}\NormalTok{ [h\_\{t}\OperatorTok{{-}}\DecValTok{1}\NormalTok{\}, x\_t])   }\CommentTok{\# update gate}
\NormalTok{r\_t }\OperatorTok{=}\NormalTok{ sigmoid(W\_r }\OperatorTok{@}\NormalTok{ [h\_\{t}\OperatorTok{{-}}\DecValTok{1}\NormalTok{\}, x\_t])   }\CommentTok{\# reset gate}

\NormalTok{h\_hat\_t }\OperatorTok{=}\NormalTok{ tanh(W }\OperatorTok{@}\NormalTok{ [r\_t }\OperatorTok{*}\NormalTok{ h\_\{t}\OperatorTok{{-}}\DecValTok{1}\NormalTok{\}, x\_t])  }\CommentTok{\# candidate}

\NormalTok{h\_t }\OperatorTok{=}\NormalTok{ (}\DecValTok{1} \OperatorTok{{-}}\NormalTok{ z\_t) }\OperatorTok{*}\NormalTok{ h\_\{t}\OperatorTok{{-}}\DecValTok{1}\NormalTok{\} }\OperatorTok{+}\NormalTok{ z\_t }\OperatorTok{*}\NormalTok{ h\_hat\_t  }\CommentTok{\# final hidden state}
\end{Highlighting}
\end{Shaded}

\begin{itemize}
\tightlist
\item
  Merges cell state and hidden state into single \texttt{h\_t}
\item
  Fewer parameters \(\rightarrow\) faster to train
\item
  Comparable performance to LSTM in practice
\end{itemize}

\subsubsection{RNN vs LSTM vs GRU}\label{rnn-vs-lstm-vs-gru}

\begin{longtable}[]{@{}
  >{\raggedright\arraybackslash}p{(\columnwidth - 6\tabcolsep) * \real{0.2379}}
  >{\raggedright\arraybackslash}p{(\columnwidth - 6\tabcolsep) * \real{0.1672}}
  >{\raggedright\arraybackslash}p{(\columnwidth - 6\tabcolsep) * \real{0.1851}}
  >{\raggedright\arraybackslash}p{(\columnwidth - 6\tabcolsep) * \real{0.4098}}@{}}
\toprule\noalign{}
\rowcolor{acc!16}
\begin{minipage}[b]{\linewidth}\raggedright
\end{minipage} & \begin{minipage}[b]{\linewidth}\raggedright
Vanilla RNN
\end{minipage} & \begin{minipage}[b]{\linewidth}\raggedright
LSTM
\end{minipage} & \begin{minipage}[b]{\linewidth}\raggedright
GRU
\end{minipage} \\
\midrule\noalign{}
\endhead
\bottomrule\noalign{}
\endlastfoot
Parameters & Fewest & Most & Middle \\
\rowcolor{zebra}
Memory length & Short & Long & Long \\
Training stability & Poor & Good & Good \\
\rowcolor{zebra}
Speed & Fastest & Slowest & Middle \\
Gates & 0 & 3 (i,f,o) & 2 (z,r) \\
\rowcolor{zebra}
When to use & Rarely & Long sequences & Long sequences, faster training \\
\end{longtable}

\subsection{Embeddings \& Word2Vec}\label{embeddings--word2vec}

\textbf{Embedding:} A dense, low-dimensional vector representation of a discrete entity (word, user, product).

\textbf{Why not one-hot?} For vocabulary of 100K words, one-hot is 100K-dimensional, sparse, and captures no semantic similarity.

\textbf{Word2Vec intuition:}

\begin{itemize}
\tightlist
\item
  Train a neural network to predict: given a word, predict its context (Skip-gram) or given context, predict center word (CBOW)
\item
  The intermediate representation (the weight matrix to the hidden layer) = word embeddings
\item
  After training: semantically similar words have similar vectors
\item
  Famous example: \texttt{king\ -\ man\ +\ woman\ ≈\ queen} (vector arithmetic captures analogy)
\end{itemize}

\textbf{Two architectures:}

\begin{itemize}
\tightlist
\item
  \textbf{CBOW (Continuous Bag of Words):} Context \(\rightarrow\) center word. Faster, good for frequent words.
\item
  \textbf{Skip-gram:} Center \(\rightarrow\) context words. Slower but better for rare words.
\end{itemize}

\textbf{Modern embeddings:} Word2Vec is conceptual predecessor to contextual embeddings (ELMo, BERT) --- same word gets different embedding depending on context (captures polysemy).

\textbf{Interview takeaway:} An embedding layer is just a lookup table (a matrix of size \texttt{vocab\_size\ ×\ embed\_dim}) --- \texttt{nn.Embedding(vocab\_size,\ d)} in PyTorch is literally a matrix, indexed by integer word IDs.

\subsection{Transfer Learning \& Fine-Tuning}\label{transfer-learning--fine-tuning}

\textbf{Transfer learning:} Use a model pre-trained on a large dataset (ImageNet, web text) as starting point for a different (usually smaller) task.

\textbf{Why it works:} Early layers learn universal features (edges, curves, basic patterns) that transfer across tasks. Later layers learn task-specific features that can be replaced or fine-tuned.

\subsubsection{Strategies}\label{strategies}

\begin{longtable}[]{@{}
  >{\raggedright\arraybackslash}p{(\columnwidth - 6\tabcolsep) * \real{0.2079}}
  >{\raggedright\arraybackslash}p{(\columnwidth - 6\tabcolsep) * \real{0.2178}}
  >{\raggedright\arraybackslash}p{(\columnwidth - 6\tabcolsep) * \real{0.1881}}
  >{\raggedright\arraybackslash}p{(\columnwidth - 6\tabcolsep) * \real{0.3861}}@{}}
\toprule\noalign{}
\rowcolor{acc!16}
\begin{minipage}[b]{\linewidth}\raggedright
Strategy
\end{minipage} & \begin{minipage}[b]{\linewidth}\raggedright
What\textquotesingle s Frozen
\end{minipage} & \begin{minipage}[b]{\linewidth}\raggedright
What\textquotesingle s Trained
\end{minipage} & \begin{minipage}[b]{\linewidth}\raggedright
When to Use
\end{minipage} \\
\midrule\noalign{}
\endhead
\bottomrule\noalign{}
\endlastfoot
Feature extraction & All pre-trained layers & Only new head & Small dataset, very similar domain \\
\rowcolor{zebra}
Fine-tuning (shallow) & Early layers & Later layers + head & Medium dataset, somewhat similar domain \\
Full fine-tuning & Nothing & Entire network & Large dataset or very different domain \\
\rowcolor{zebra}
LoRA/PEFT & Most weights & Low-rank adapters & LLM fine-tuning, minimal compute \\
\end{longtable}

\textbf{Key concepts:}

\begin{itemize}
\tightlist
\item
  \textbf{Pre-trained model:} Trained on large, general dataset (ImageNet for vision, BookCorpus/Wikipedia for NLP)
\item
  \textbf{Downstream task:} Your specific task (medical image classification, sentiment analysis)
\item
  \textbf{Head:} The final layer(s) replaced for the new task (e.g., new linear classifier)
\item
  \textbf{Learning rate:} Use much smaller LR for pre-trained layers (1e-5 to 1e-4) vs new layers (1e-3)
\end{itemize}

\textbf{Interview takeaway:} Transfer learning is one of the most practically important concepts in modern ML --- almost no one trains from scratch at FAANG; everything starts from a pre-trained checkpoint.

\subsection{Autoencoders}\label{autoencoders}

An \textbf{autoencoder} is a neural network trained to compress its input into a low-dimensional \textbf{latent representation} (encoding) and then reconstruct the original input (decoding):

\setcodelang{}

\begin{verbatim}
Input x → Encoder → Latent z (bottleneck) → Decoder → Reconstruction x̂

Loss = ||x - x̂||²  (reconstruction loss)
\end{verbatim}

\textbf{Why it\textquotesingle s useful:}

\begin{itemize}
\tightlist
\item
  Learns unsupervised feature representations
\item
  Anomaly detection: reconstruction error high for anomalous inputs
\item
  Dimensionality reduction (non-linear PCA)
\item
  \textbf{Variational Autoencoder (VAE):} Latent space is a probability distribution (Normal) \(\rightarrow\) enables generation. Loss adds KL-divergence term.
\end{itemize}

\subsection{Generative Adversarial Networks (GANs)}\label{generative-adversarial-networks-gans}

\textbf{Two networks competing in a minimax game:}

\setcodelang{}

\begin{verbatim}
Generator G: latent noise z → fake sample G(z)
Discriminator D: real/fake sample → probability of being real

G tries to fool D: maximize log(D(G(z)))
D tries to catch G: maximize log(D(x)) + log(1 - D(G(z)))

Combined: min_G max_D V(D,G) = E[log D(x)] + E[log(1-D(G(z)))]
\end{verbatim}

\textbf{Training loop:}

\begin{enumerate}
\def\labelenumi{\arabic{enumi}.}
\tightlist
\item
  Sample real data batch \texttt{x}
\item
  Sample noise \texttt{z}, generate fake data \texttt{G(z)}
\item
  Train D to distinguish real from fake
\item
  Train G to fool D (freeze D\textquotesingle s weights)
\item
  Repeat
\end{enumerate}

\textbf{Key challenges:}

\begin{itemize}
\tightlist
\item
  \textbf{Mode collapse:} G learns to generate only a few modes (types) of data
\item
  \textbf{Training instability:} G and D need to progress together
\item
  \textbf{Evaluation:} FID (Fréchet Inception Distance) score
\end{itemize}

\textbf{Modern successors:} Diffusion models have largely replaced GANs for image generation (DALL-E, Stable Diffusion) due to training stability.

\subsection{Overfitting in Deep Learning}\label{overfitting-in-deep-learning}

\textbf{Overfitting:} Model memorizes training data, performs poorly on unseen data. Neural nets are especially prone --- they have massive capacity.

\textbf{Diagnosing:}

\setcodelang{}

\begin{verbatim}
Training loss: 0.01
Validation loss: 2.5  → Classic overfit
\end{verbatim}

\textbf{Solutions in DL:}

\begin{longtable}[]{@{}
  >{\raggedright\arraybackslash}p{(\columnwidth - 2\tabcolsep) * \real{0.3120}}
  >{\raggedright\arraybackslash}p{(\columnwidth - 2\tabcolsep) * \real{0.6880}}@{}}
\toprule\noalign{}
\rowcolor{acc!16}
\begin{minipage}[b]{\linewidth}\raggedright
Technique
\end{minipage} & \begin{minipage}[b]{\linewidth}\raggedright
Effect
\end{minipage} \\
\midrule\noalign{}
\endhead
\bottomrule\noalign{}
\endlastfoot
More data & Best solution; data augmentation as proxy \\
\rowcolor{zebra}
Dropout & Ensemble effect, breaks co-adaptation \\
Weight decay (L2) & Smaller weights \(\rightarrow\) smoother function \\
\rowcolor{zebra}
Early stopping & Stop before model memorizes training set \\
Batch normalization & Mild regularization via batch noise \\
\rowcolor{zebra}
Reduce model capacity & Fewer layers/neurons \\
Data augmentation & Random transformations \(\rightarrow\) effective more data \\
\rowcolor{zebra}
Label smoothing & Prevents overconfident predictions \\
\end{longtable}

\textbf{Bias-variance in DL:}

\begin{itemize}
\tightlist
\item
  Underfitting = high bias \(\rightarrow\) need bigger model, more training
\item
  Overfitting = high variance \(\rightarrow\) need regularization, more data
\item
  Deep learning has largely killed the classical bias-variance trade-off concern at scale --- modern large models can simultaneously have low bias AND low variance if trained with enough data (the "double descent" phenomenon)
\end{itemize}

\section{Architecture / Diagrams}\label{architecture--diagrams}

\subsection{Single Neuron / Perceptron}\label{single-neuron--perceptron}

\setcodelang{JavaScript}

\begin{Shaded}
\begin{Highlighting}[]
\NormalTok{    x1 }\OperatorTok{{-}{-}{-}{-}}\NormalTok{w1}\OperatorTok{{-}{-}{-}{-}}\NormalTok{\textbackslash{}}
\NormalTok{                  \textbackslash{}}
\NormalTok{    x2 }\OperatorTok{{-}{-}{-}{-}}\NormalTok{w2}\OperatorTok{{-}{-}{-}{-}{-}}\NormalTok{[∑ }\OperatorTok{+}\NormalTok{ b]}\OperatorTok{{-}{-}{-}{-}}\NormalTok{[}\FunctionTok{f}\NormalTok{(}\OperatorTok{.}\NormalTok{)]}\OperatorTok{{-}{-}{-}{-}{-}\textgreater{}}\NormalTok{ y}
                  \OperatorTok{/}
\NormalTok{    x3 }\OperatorTok{{-}{-}{-}{-}}\NormalTok{w3}\OperatorTok{{-}{-}{-}{-}}\SpecialStringTok{/}

\NormalTok{∑ }\OperatorTok{=}\NormalTok{ w1}\OperatorTok{*}\NormalTok{x1 }\OperatorTok{+}\NormalTok{ w2}\OperatorTok{*}\NormalTok{x2 }\OperatorTok{+}\NormalTok{ w3}\OperatorTok{*}\NormalTok{x3 }\OperatorTok{+} \FunctionTok{b}\NormalTok{   (linear combination)}
\NormalTok{f }\OperatorTok{=}\NormalTok{ activation }\KeywordTok{function}\NormalTok{           (non}\OperatorTok{{-}}\NormalTok{linearity)}
\NormalTok{y }\OperatorTok{=} \FunctionTok{f}\NormalTok{(∑)                          (neuron output)}
\end{Highlighting}
\end{Shaded}

\subsection{MLP Architecture}\label{mlp-architecture}

\visualfigure{../figures/aiml-03-dl/03-mlp-network.pdf}{A feed-forward net stacks fully-connected layers: each neuron is a weighted sum of the previous layer plus a bias, passed through a nonlinearity.}

\subsection{Forward + Backpropagation Flow}\label{forward--backpropagation-flow}

\setcodelang{}

\begin{verbatim}
FORWARD PASS (left to right, computing predictions):
────────────────────────────────────────────────────────
x ──→ [Layer 1: z¹=W¹x+b¹, a¹=f(z¹)] ──→ [Layer 2: z²=W²a¹+b², a²=f(z²)] ──→ ŷ
                                                                                    │
                                                                                [Loss L]
                                                                                    │
BACKWARD PASS (right to left, computing gradients):                                 ↓
────────────────────────────────────────────────────────                       dL/dŷ
                                                                                    │
←── [dL/dW¹ = dL/da² · da²/dz² · dz²/dW²] ←── [dL/da² = dL/dŷ · dŷ/da²] ←──────
        │                                                       │
   Update W²                                              chain rule
   
Chain Rule:
dL    dL   da²   dz²
── = ─── · ─── · ───
dW²  da²   dz²  dW²

    = dL/da² · f'(z²) · a¹ᵀ
\end{verbatim}

\subsection{Activation Function Shapes}\label{activation-function-shapes}

\setcodelang{}

\begin{verbatim}
SIGMOID:            TANH:               ReLU:               Leaky ReLU:
    1 ┤  ___         1 ┤  ___              │    /            │   /
      │ /               │ /               │   /             │  /
  0.5 ┤/            0   ┤                 │  /          0   ┤ /
      │              -1 ┤___              │/                │/
  0 ──┤──────           ─────────     ────┤────         ────┤────
     -∞  +∞            -∞    +∞         0│              -0.01\
                                                              \

Saturates at ±∞   Saturates at ±∞   Dead for x<0    Small slope for x<0
\end{verbatim}

\subsection{CNN Architecture}\label{cnn-architecture}

\visualfigure{../figures/aiml-03-dl/04-cnn-architecture.pdf}{Convolution + pooling progressively reduce spatial size while increasing channel depth, building features from edges $\rightarrow$ textures $\rightarrow$ objects, then a dense head classifies.}

\subsection{Gradient Descent Visualization}\label{gradient-descent-visualization}

\setcodelang{}

\begin{verbatim}
Loss
  │
  │ \
  │  \         Learning rate too large: overshoots
  │   \  X X X
  │    \ /
  │     *        Optimal: converges smoothly
  │    / \
  │   /   \
  └──────────── Weights (W)
  
  θ_{t+1} = θ_t - α * ∇L(θ_t)
  
  Too large α → diverge or oscillate
  Too small α → very slow convergence
  Adam → adaptive α per parameter
\end{verbatim}

\subsection{RNN / LSTM Unrolled}\label{rnn--lstm-unrolled}

\setcodelang{}

\begin{verbatim}
VANILLA RNN (unrolled across time):

h0──→[RNN]──h1──→[RNN]──h2──→[RNN]──h3
      ↑              ↑              ↑
      x1             x2             x3
      
Same weights W_xh, W_hh at every step.

LSTM CELL (single time step):

        h_{t-1} ──────────────────────────────→
                    │        │        │        │
        x_t ────────┼────────┼────────┼────────┼──→
                    ↓        ↓        ↓        ↓
               [Forget  [Input   [Cell    [Output
                Gate f]  Gate i]  Gate g]  Gate o]
               sigmoid  sigmoid   tanh   sigmoid
                    │        │        │        │
c_{t-1}─────→[  × f_t ]──→[+ i_t*g_t]──→[tanh]─→[× o_t]──→ h_t
                                    c_t                        ↓
                                                            output
                                                            
KEY: c_t uses ADDITION (+), not multiplication → gradient highway!
\end{verbatim}

\section{Real-World Examples}\label{real-world-examples}

\begin{longtable}[]{@{}
  >{\raggedright\arraybackslash}p{(\columnwidth - 6\tabcolsep) * \real{0.1287}}
  >{\raggedright\arraybackslash}p{(\columnwidth - 6\tabcolsep) * \real{0.2673}}
  >{\raggedright\arraybackslash}p{(\columnwidth - 6\tabcolsep) * \real{0.2970}}
  >{\raggedright\arraybackslash}p{(\columnwidth - 6\tabcolsep) * \real{0.3069}}@{}}
\toprule\noalign{}
\rowcolor{acc!16}
\begin{minipage}[b]{\linewidth}\raggedright
Company
\end{minipage} & \begin{minipage}[b]{\linewidth}\raggedright
Application
\end{minipage} & \begin{minipage}[b]{\linewidth}\raggedright
Architecture
\end{minipage} & \begin{minipage}[b]{\linewidth}\raggedright
Key Technique
\end{minipage} \\
\midrule\noalign{}
\endhead
\bottomrule\noalign{}
\endlastfoot
Google Photos & Image search/classification & EfficientNet/ViT & CNN pre-trained on JFT-300M \\
\rowcolor{zebra}
Netflix & Thumbnail CTR optimization & CNN & Transfer from ImageNet \\
Amazon Alexa & Wake word detection & Small CNN/RNN & Trained to run on device \\
\rowcolor{zebra}
YouTube & Video recommendation & Two-Tower DNN & Embeddings for users and videos \\
Facebook & Hate speech detection & BERT fine-tuned & Transfer learning \\
\rowcolor{zebra}
Tesla & Autopilot vision & Multi-camera CNN & Custom hardware (FSD chip) \\
Spotify & Song recommendations & Autoencoder + collab filtering & Learned audio embeddings \\
\rowcolor{zebra}
OpenAI DALL-E & Image generation & Diffusion + CLIP & Contrastive + generative \\
\end{longtable}

\section{Real-Life Analogies --- The Learning Factory Assembly Line}\label{real-life-analogies--the-learning-factory-assembly-line}

\emph{Every component of deep learning maps onto a factory that manufactures products and learns by inspecting its own output quality.}

\begin{longtable}[]{@{}
  >{\raggedright\arraybackslash}p{(\columnwidth - 2\tabcolsep) * \real{0.1775}}
  >{\raggedright\arraybackslash}p{(\columnwidth - 2\tabcolsep) * \real{0.8225}}@{}}
\toprule\noalign{}
\rowcolor{acc!16}
\begin{minipage}[b]{\linewidth}\raggedright
DL Concept
\end{minipage} & \begin{minipage}[b]{\linewidth}\raggedright
Factory Analogy
\end{minipage} \\
\midrule\noalign{}
\endhead
\bottomrule\noalign{}
\endlastfoot
Input data & Raw materials entering the factory \\
\rowcolor{zebra}
Neural network & The entire factory assembly line \\
Layer & One station on the assembly line (e.g., cutting, welding, painting) \\
\rowcolor{zebra}
Weights & The dial settings at each station (speed, temperature, pressure) \\
Forward pass & The product traveling down the line, being transformed at each station \\
\rowcolor{zebra}
Activation function & A station\textquotesingle s go/no-go quality gate --- product passes or gets modified \\
Loss function & The final Quality Control (QC) inspector\textquotesingle s defect score \\
\rowcolor{zebra}
Backpropagation & The foreman walking the line \textbf{backward} from QC, telling each station specifically how its settings caused the defect \\
Gradient & How much each station\textquotesingle s dials need to change to reduce defects \\
\rowcolor{zebra}
Gradient descent & The foreman nudging each dial a small amount in the right direction \\
Learning rate & How big a nudge the foreman applies (too big = overcompensate; too small = too slow) \\
\rowcolor{zebra}
Adam optimizer & A smart foreman who tracks which dials have been moved a lot recently and adjusts nudge size per-dial accordingly \\
Vanishing gradient & By the time the foreman reaches station 1 (near the start), his feedback has been diluted so many times he can only whisper --- station 1 barely learns \\
\rowcolor{zebra}
Residual connections & A direct phone line from QC to every station, so even station 1 gets clear feedback \\
Dropout & Randomly sending workers home each day so the line can\textquotesingle t over-rely on any one person --- forces robustness \\
\rowcolor{zebra}
Batch normalization & Standardizing the product dimensions at each station so later stations don\textquotesingle t have to compensate for upstream variation \\
CNN & Inspectors using a small magnifying glass that slides across the product surface, looking for defects in each patch --- same inspector logic applied everywhere \\
\rowcolor{zebra}
RNN & A line that remembers the last product it processed, using that memory to better process the current one \\
LSTM & A line with a long-term clipboard (cell state) that records important product history, with explicit "write/erase/read" decisions \\
\rowcolor{zebra}
Transfer learning & Hiring an experienced foreman from another factory (pre-trained on a related product) --- retrain only the final stations for your specific product \\
Overfitting & The line gets so optimized for the specific training products it\textquotesingle s seen that it fails on new, slightly different products \\
\rowcolor{zebra}
Embeddings & Each product type gets a detailed spec sheet (dense vector) summarizing its properties for downstream stations \\
\end{longtable}

\section{Memory Tricks / Mnemonics}\label{memory-tricks--mnemonics}

\textbf{Activation functions --- "STRL GS":}

\begin{itemize}
\tightlist
\item
  \textbf{S}igmoid: 0-1 output, squashes, \textbf{S}aturates
\item
  \textbf{T}anh: -1 to 1, \textbf{T}wo-sided sigmoid, zero-centered
\item
  \textbf{R}eLU: \textbf{R}ectify negatives to zero, \textbf{R}isk of dead neurons
\item
  \textbf{L}eaky ReLU: \textbf{L}ittle slope for negatives
\item
  \textbf{G}ELU: \textbf{G}aussian-flavored, used in \textbf{G}PT
\item
  \textbf{S}oftmax: \textbf{S}um to 1, for \textbf{S}election (multiclass)
\end{itemize}

\textbf{Optimizer progression --- "S M R A":}

\begin{itemize}
\tightlist
\item
  \textbf{S}GD \(\rightarrow\) add \textbf{M}omentum \(\rightarrow\) add per-param LR (\textbf{R}MSprop) \(\rightarrow\) combine both (\textbf{A}dam)
\end{itemize}

\textbf{LSTM gates --- "I F O" (IFO = "I Forget Output"):}

\begin{itemize}
\tightlist
\item
  \textbf{I}nput gate: what new info to write to cell
\item
  \textbf{F}orget gate: what old info to erase from cell
\item
  \textbf{O}utput gate: what to read from cell as hidden state
\end{itemize}

\textbf{Weight init --- "Xavier for tanh, He for ReLU":}

\begin{itemize}
\tightlist
\item
  Xavier \(\rightarrow\) "X" looks like a cross \(\rightarrow\) tanh crosses zero
\item
  He \(\rightarrow\) "H" \(\rightarrow\) Half inputs are zero (ReLU), so scale by 2/n
\end{itemize}

\textbf{BatchNorm vs LayerNorm --- "Batch = across samples, Layer = across features":}

\begin{itemize}
\tightlist
\item
  \textbf{B}atch: \textbf{B}ig dimension (across many samples per feature)
\item
  \textbf{L}ayer: \textbf{L}ength of the feature vector (across all features per sample)
\end{itemize}

\textbf{CNN parameter count formula --- "K K C + 1 times F":}

\begin{itemize}
\tightlist
\item
  (Kernel\_h \(\times\) Kernel\_w \(\times\) Channels\_in + 1) \(\times\) Filters\_out
\end{itemize}

\textbf{Backprop chain rule --- "LOCAL \(\times\) UPSTREAM":}

\begin{itemize}
\tightlist
\item
  At each node: gradient = local derivative \(\times\) gradient from downstream
\item
  \texttt{dL/dW\ =\ (dL/da\_out)\ ×\ (da\_out/dz)\ ×\ (dz/dW)}
\item
  = upstream gradient \(\times\) activation derivative \(\times\) input
\end{itemize}

\textbf{Vanishing gradient solutions --- "RICH":}

\begin{itemize}
\tightlist
\item
  \textbf{R}eLU (better activation)
\item
  \textbf{I}nitialization (Xavier/He)
\item
  \textbf{C}onnections (residual/skip)
\item
  \textbf{H}normalization (batch/layer norm)
\end{itemize}

\section{Common Interview Questions}\label{common-interview-questions}

\subsection{Q1: Explain backpropagation from first principles.}\label{q1-explain-backpropagation-from-first-principles}

\textbf{Model answer:}
Backpropagation computes gradients of the loss with respect to all parameters using the chain rule of calculus. After the forward pass computes the loss, we propagate the error signal backward through the computational graph.

At each layer, we compute three things:

\begin{enumerate}
\def\labelenumi{\arabic{enumi}.}
\tightlist
\item
  \texttt{dL/dz\^{}l\ =\ dL/da\^{}l\ *\ f\textquotesingle{}(z\^{}l)} --- gradient through activation
\item
  \texttt{dL/dW\^{}l\ =\ a\^{}(l-1)\^{}T\ *\ dL/dz\^{}l} --- gradient for weight update
\item
  \texttt{dL/da\^{}(l-1)\ =\ dL/dz\^{}l\ *\ W\^{}l\^{}T} --- gradient to pass to previous layer
\end{enumerate}

The key insight: each layer\textquotesingle s gradient depends only on the next layer\textquotesingle s gradient (upstream) and its own local derivatives. We store forward pass activations and reuse them in the backward pass.

\textbf{Follow-up:} "Why do we need to store activations from the forward pass?"

\begin{itemize}
\tightlist
\item
  We need \texttt{a\^{}(l-1)} to compute \texttt{dL/dW\^{}l\ =\ a\^{}(l-1)\^{}T\ *\ dL/dz\^{}l}
\item
  We need \texttt{z\^{}l} to compute \texttt{f\textquotesingle{}(z\^{}l)} (the activation derivative)
\item
  Memory cost of training = O(num\_layers \(\times\) batch\_size \(\times\) layer\_size) for stored activations
\item
  Gradient checkpointing trades compute for memory by recomputing activations
\end{itemize}

\subsection{Q2: Why does ReLU work better than sigmoid for hidden layers?}\label{q2-why-does-relu-work-better-than-sigmoid-for-hidden-layers}

\textbf{Model answer:}
Three reasons:

\begin{enumerate}
\def\labelenumi{\arabic{enumi}.}
\item
  \textbf{No vanishing gradient in the positive region.} Sigmoid derivative max is 0.25 --- multiply that through 50 layers and the gradient \(\approx\) 0. ReLU derivative for x\textgreater0 is exactly 1, so gradients flow without shrinking.
\item
  \textbf{Computationally cheap.} ReLU is just \texttt{max(0,\ x)} --- a single comparison. Sigmoid requires an exponentiation.
\item
  \textbf{Sparse activation.} Only \textasciitilde50\% of neurons activate for any given input. Sparse representations are more computationally efficient and have been linked to better generalization.
\end{enumerate}

\textbf{Trade-off:} Dead neurons --- if a neuron\textquotesingle s weights cause it to always receive negative inputs, gradient = 0 forever, weights never update. Solutions: Leaky ReLU, careful initialization, lower learning rate.

\textbf{Follow-up:} "When would you still use sigmoid?"

\begin{itemize}
\tightlist
\item
  Output layer for binary classification (need probability in {[}0,1{]})
\item
  LSTM gates (need to produce values in {[}0,1{]} for gating)
\end{itemize}

\subsection{Q3: Explain the vanishing gradient problem and how ResNets solve it.}\label{q3-explain-the-vanishing-gradient-problem-and-how-resnets-solve-it}

\textbf{Model answer:}
In backprop, the gradient at layer \texttt{l} is the product of Jacobians from layer \texttt{l} to the output. With L layers of sigmoid (max derivative 0.25), the gradient magnitude scales as 0.25\^{}L. For L=20, that\textquotesingle s \textasciitilde10\^{}-12 --- effectively zero. Early layers learn nothing.

ResNets add skip connections: \texttt{output\ =\ F(x,\ W)\ +\ x}

During backprop: \texttt{dL/dx\ =\ dL/d(output)\ *\ (dF/dx\ +\ I)}

The \texttt{I} (identity) term means there\textquotesingle s always a direct gradient path through the skip connection --- gradient = upstream + local. Even if \texttt{dF/dx} is small, the upstream gradient flows unimpeded. This lets ResNets be 50-152 layers deep without gradient issues.

\textbf{Analogy:} Like adding a fiber optic cable directly from the foreman\textquotesingle s office to every station, bypassing the chain of verbal relays.

\subsection{Q4: Compare Adam vs SGD. When would you use each?}\label{q4-compare-adam-vs-sgd-when-would-you-use-each}

\textbf{Model answer:}

\textbf{Adam advantages:}

\begin{itemize}
\tightlist
\item
  Adaptive per-parameter learning rates --- features with rare but informative gradients (like word embeddings for rare words) get appropriate updates
\item
  Incorporates momentum automatically
\item
  Requires less tuning --- default hyperparameters work well in most cases
\item
  Converges faster in wall-clock time
\end{itemize}

\textbf{SGD+Momentum advantages:}

\begin{itemize}
\tightlist
\item
  Often generalizes slightly better (converges to flatter minima)
\item
  Lower memory usage (2x params vs 3x for Adam)
\item
  Better understood theoretically
\end{itemize}

\textbf{When to use each:}

\begin{itemize}
\tightlist
\item
  Adam: Default choice, especially for: NLP tasks, transformers, when prototyping, when data is sparse
\item
  SGD+Momentum: When ultimate accuracy matters and you can afford more tuning (ResNet image classifiers), fine-tuning of pre-trained models
\end{itemize}

\textbf{AdamW:} Adam with proper weight decay (decoupled from gradient update). Use this for transformer fine-tuning --- standard Adam\textquotesingle s weight decay is mathematically incorrect (it interacts with the adaptive LR).

\subsection{Q5: Why does dropout work as regularization?}\label{q5-why-does-dropout-work-as-regularization}

\textbf{Model answer:}
Two complementary explanations:

\begin{enumerate}
\def\labelenumi{\arabic{enumi}.}
\item
  \textbf{Ensemble interpretation:} With p=0.5 dropout on n neurons, there are 2\^{}n possible sub-networks. Training with dropout trains all of them simultaneously with shared weights. At inference, the full network (with weights scaled by 1-p) approximates the average of this ensemble.
\item
  \textbf{Co-adaptation prevention:} Without dropout, neurons can develop complex codependencies --- Neuron A learns to always correct Neuron B\textquotesingle s mistakes. This is fragile. Dropout forces each neuron to be useful independently, learning more robust features.
\end{enumerate}

\textbf{Implementation detail:} During training, multiply by mask and scale by 1/(1-p). During inference, no dropout (or equivalently, scale weights by (1-p)). This is called "inverted dropout" --- keeps expected values consistent between train and inference.

\textbf{Follow-up:} "Why not use dropout in batch normalization layers or right before softmax?"

\begin{itemize}
\tightlist
\item
  BatchNorm uses batch statistics that become noisy with dropout
\item
  Dropout after softmax would create non-probability outputs
\end{itemize}

\subsection{Q6: What\textquotesingle s the difference between batch normalization and layer normalization?}\label{q6-whats-the-difference-between-batch-normalization-and-layer-normalization}

\textbf{Model answer:}
Both normalize activations to stabilize training, but they differ in \textbf{which dimension} they normalize over:

\textbf{BatchNorm:} Normalizes across the batch dimension, per feature.

\begin{itemize}
\tightlist
\item
  For a batch of B samples, each with N features: compute mean/var across B for each of the N features
\item
  Problem: behavior depends on batch size; different at train vs inference (uses running stats)
\item
  Great for: CNNs (images), large batch sizes
\end{itemize}

\textbf{LayerNorm:} Normalizes across the feature dimension, per sample.

\begin{itemize}
\tightlist
\item
  For each sample independently: compute mean/var across its N features
\item
  Same behavior at train and inference
\item
  Works for any batch size (including batch\_size=1)
\item
  Required for: Transformers, RNNs, variable-length sequences
\end{itemize}

\textbf{Intuition:} BatchNorm says "make this feature zero-mean unit-variance across the batch." LayerNorm says "make this sample\textquotesingle s representation zero-mean unit-variance across features."

\subsection{Q7: How do LSTMs solve the vanishing gradient problem in RNNs?}\label{q7-how-do-lstms-solve-the-vanishing-gradient-problem-in-rnns}

\textbf{Model answer:}
Vanilla RNNs: \texttt{h\_t\ =\ tanh(W\_h\ *\ h\_\{t-1\}\ +\ W\_x\ *\ x\_t)}. The gradient through time requires multiplying \texttt{dh\_t/dh\_\{t-1\}\ =\ W\_h\ *\ tanh\textquotesingle{}(...)} for T steps \(\rightarrow\) exponentially small.

LSTMs use an \textbf{additive cell state update:}
\texttt{c\_t\ =\ f\_t\ *\ c\_\{t-1\}\ +\ i\_t\ *\ g\_t}

The \texttt{+} (addition) is the key. The gradient of \texttt{c\_t} w.r.t. \texttt{c\_\{t-1\}} includes the forget gate \texttt{f\_t}. When \texttt{f\_t\ ≈\ 1} (remember everything), the gradient flows unimpeded: \texttt{dc\_t/dc\_\{t-1\}\ =\ f\_t\ ≈\ 1}. This is the same "gradient highway" trick as ResNets.

The LSTM can learn to keep its forget gate open for long-range dependencies, allowing gradients to flow back hundreds of steps.

\textbf{Follow-up:} "Do LSTMs completely solve vanishing gradients?"
No --- they mitigate it significantly but don\textquotesingle t eliminate it. For very long sequences (\textgreater500 steps), attention mechanisms (Transformers) work better because every position can directly attend to every other position.

\subsection{Q8: Explain transfer learning and when fine-tuning is appropriate.}\label{q8-explain-transfer-learning-and-when-fine-tuning-is-appropriate}

\textbf{Model answer:}
Transfer learning uses knowledge from a pre-trained model (trained on a large dataset) as the starting point for a new task. Instead of random initialization, start from a model that already understands useful features.

\textbf{When to use which strategy:}

\begin{longtable}[]{@{}
  >{\raggedright\arraybackslash}p{(\columnwidth - 4\tabcolsep) * \real{0.2009}}
  >{\raggedright\arraybackslash}p{(\columnwidth - 4\tabcolsep) * \real{0.2439}}
  >{\raggedright\arraybackslash}p{(\columnwidth - 4\tabcolsep) * \real{0.5552}}@{}}
\toprule\noalign{}
\rowcolor{acc!16}
\begin{minipage}[b]{\linewidth}\raggedright
Your Data Size
\end{minipage} & \begin{minipage}[b]{\linewidth}\raggedright
Domain Similarity
\end{minipage} & \begin{minipage}[b]{\linewidth}\raggedright
Strategy
\end{minipage} \\
\midrule\noalign{}
\endhead
\bottomrule\noalign{}
\endlastfoot
Small & Similar & Freeze all, train only head \\
\rowcolor{zebra}
Small & Different & Freeze early layers, fine-tune later layers + head \\
Large & Similar & Fine-tune entire network with small LR \\
\rowcolor{zebra}
Large & Different & Full fine-tuning or train from scratch \\
\end{longtable}

\textbf{Practical tips:}

\begin{itemize}
\tightlist
\item
  Use learning rate \textasciitilde10-100x smaller for pre-trained layers than new layers
\item
  Fine-tune in stages (unfreeze more layers progressively)
\item
  Warm up learning rate when unfreezing large pre-trained models
\end{itemize}

\textbf{Why it works:} Early CNN layers detect edges and textures that are universal across vision tasks. Early NLP layers learn syntax/semantics that transfer across text tasks. Only the final, task-specific layers need heavy adaptation.

\section{Senior-Level Discussion Points}\label{senior-level-discussion-points}

\subsection{1. The Double Descent Phenomenon}\label{1-the-double-descent-phenomenon}

Classical statistics says: more model capacity \(\rightarrow\) overfitting. Modern deep learning violates this. As model capacity grows beyond the interpolation threshold (where training loss = 0), test loss can \emph{decrease again} --- the "second descent." Large models that perfectly fit training data can still generalize well. This challenges bias-variance trade-off intuitions.

\subsection{2. Neural Scaling Laws}\label{2-neural-scaling-laws}

Performance improves predictably as a power law with: number of parameters, dataset size, and compute. Kaplan et al. (OpenAI, 2020) showed these laws hold over many orders of magnitude. \textbf{Implication:} Given a compute budget, optimal allocation is a specific ratio of model size to data (Chinchilla scaling).

\subsection{3. Loss Landscape Geometry}\label{3-loss-landscape-geometry}

SGD with noise tends to find flatter minima (wider loss basins) than Adam. Flat minima generalize better --- small perturbations to weights cause small changes in loss. Sharp minima are "brittle." This is one reason SGD sometimes generalizes better than Adam despite slower convergence.

\subsection{4. Gradient Checkpointing}\label{4-gradient-checkpointing}

Trading compute for memory: instead of storing all activations during forward pass, recompute them during backward pass. Reduces memory from O(L) to O(\(\surd\)L) layers at the cost of \textasciitilde33\% more compute. Essential for training very deep models on limited GPU memory.

\subsection{5. Mixed Precision Training}\label{5-mixed-precision-training}

Use FP16 for forward/backward pass, FP32 for weight updates. 2x memory savings, \textasciitilde3x throughput on modern GPUs. Critical: maintain FP32 master copy of weights, use loss scaling to prevent FP16 underflow of small gradients.

\subsection{6. BatchNorm at Inference vs Training}\label{6-batchnorm-at-inference-vs-training}

BatchNorm uses batch statistics during training (noisy, slightly different each step) but uses running averages at inference (deterministic). This mismatch can cause subtle bugs --- if your training pipeline has any distribution shift between batches, the running statistics diverge from batch statistics. Always call \texttt{model.eval()} before inference in PyTorch.

\subsection{7. The Problem with Adam\textquotesingle s Weight Decay}\label{7-the-problem-with-adams-weight-decay}

Standard Adam applies weight decay to the adaptive-LR-scaled gradient, not the weight directly. This means weight decay strength varies by parameter history. AdamW decouples weight decay: \texttt{W\ =\ W\ *\ (1\ -\ wd)\ -\ lr\ *\ m\_hat\ /\ (sqrt(v\_hat)\ +\ ε)}. Always use AdamW for Transformers.

\subsection{8. Dying ReLU in Practice}\label{8-dying-relu-in-practice}

Networks with very large learning rates can have many neurons "die" in the first few training steps. Symptoms: training loss stops decreasing, gradient norms collapse. Solutions: reduce LR, use He initialization, use Leaky ReLU, use gradient clipping.

\section{Typical Mistakes Candidates Make}\label{typical-mistakes-candidates-make}

\begin{enumerate}
\def\labelenumi{\arabic{enumi}.}
\item
  \textbf{"Backpropagation updates weights during the forward pass"} --- No. Forward pass computes and stores activations. Backward pass computes gradients. Optimizer uses gradients to update weights. Three separate steps.
\item
  \textbf{"ReLU solves vanishing gradients completely"} --- ReLU has gradient 1 for x\textgreater0, but dead neurons (gradient=0 for x\textless0) are another failure mode. Also, stacking many ReLU layers with no skip connections can still have issues.
\item
  \textbf{"Dropout is applied the same way at train and inference"} --- No. Dropout is deactivated at inference time. In PyTorch: \texttt{model.eval()} turns off dropout; \texttt{model.train()} turns it back on.
\item
  \textbf{"BatchNorm removes the need for weight initialization"} --- BatchNorm helps a lot but good initialization still matters, especially for the first few steps before BatchNorm statistics are meaningful.
\item
  \textbf{"Softmax in hidden layers"} --- Softmax is for output layer only. Using softmax in hidden layers is almost always wrong (it creates competition between neurons where you want independence).
\item
  \textbf{"Larger learning rate always means faster training"} --- Too large \(\rightarrow\) divergence, oscillation, overshooting. Learning rate needs to be paired with batch size: if you double batch size, scale LR by \(\surd\)2 or 2 (linear scaling rule).
\item
  \textbf{"More layers always helps"} --- Without residual connections, adding layers often hurts (degradation problem). Depth needs to be matched with appropriate architectural choices.
\item
  \textbf{"Adam always converges to the best solution"} --- Adam converges faster but to different (sometimes worse generalizing) minima than SGD+momentum in some tasks. For production image classifiers (ResNets), SGD+momentum with careful LR schedule often wins.
\item
  \textbf{"Word2Vec gives the same embedding for \textquotesingle bank\textquotesingle{} (river) and \textquotesingle bank\textquotesingle{} (financial)"} --- Static embeddings (Word2Vec) have one vector per word regardless of context. Contextual embeddings (ELMo, BERT) solve this.
\item
  \textbf{Confusing parameters vs hyperparameters} --- Weights/biases are parameters (learned via gradient descent). Learning rate, batch size, number of layers, dropout rate, etc. are hyperparameters (set before training, tuned via validation).
\end{enumerate}

\section{How This Connects to Other Topics}\label{how-this-connects-to-other-topics}

\subsection{\texorpdfstring{\(\rightarrow\) ML Fundamentals}{\textbackslash rightarrow ML Fundamentals}}\label{rightarrow-ml-fundamentals}

\begin{itemize}
\tightlist
\item
  Gradient descent originated in classical ML (logistic regression, SVMs with SGD)
\item
  Bias-variance trade-off, regularization (L1/L2), cross-validation: same concepts apply
\item
  Neural networks are function approximators --- they\textquotesingle re just more powerful versions of linear models with non-linear feature transformations
\end{itemize}

\subsection{\texorpdfstring{\(\rightarrow\) Transformers / LLMs}{\textbackslash rightarrow Transformers / LLMs}}\label{rightarrow-transformers--llms}

\begin{itemize}
\tightlist
\item
  Self-attention is the core operation replacing RNN sequential processing
\item
  LayerNorm is used instead of BatchNorm (batch-size independence)
\item
  Embedding layers + positional encodings build directly on Word2Vec ideas
\item
  Pre-training + fine-tuning transfer learning at massive scale
\item
  GELU activation (vs ReLU) is standard in transformers
\end{itemize}

\subsection{\texorpdfstring{\(\rightarrow\) MLOps}{\textbackslash rightarrow MLOps}}\label{rightarrow-mlops}

\begin{itemize}
\tightlist
\item
  Training vs inference mode (BatchNorm, Dropout) must be handled correctly in serving
\item
  Model serialization: saving/loading weights (\texttt{.pt}, \texttt{.safetensors}, ONNX)
\item
  GPU memory management: gradient checkpointing, mixed precision
\item
  Distributed training: data parallelism (split batch across GPUs), model parallelism (split model across GPUs)
\item
  Experiment tracking: loss curves, gradient norms, activation distributions (TensorBoard, W\&B)
\end{itemize}

\subsection{\texorpdfstring{\(\rightarrow\) Performance Engineering}{\textbackslash rightarrow Performance Engineering}}\label{rightarrow-performance-engineering}

\begin{itemize}
\tightlist
\item
  Operator fusion: combining Conv + BatchNorm + ReLU into single kernel
\item
  Quantization: FP32 \(\rightarrow\) INT8 for inference (4x memory, 2-4x throughput)
\item
  Knowledge distillation: small model trained to mimic large model\textquotesingle s soft outputs
\item
  Pruning: remove unimportant weights (small magnitude), structured pruning for hardware efficiency
\item
  Hardware: GPU tensor cores optimized for matrix multiplication, TPUs for large batch training
\end{itemize}

\section{FAANG Interview Tips}\label{faang-interview-tips}

\textbf{Coding questions (ML-adjacent):}

\begin{itemize}
\tightlist
\item
  Know how to implement forward propagation in numpy: \texttt{h\ =\ np.maximum(0,\ X\ @\ W\ +\ b)} (ReLU + linear layer)
\item
  Know how to compute cross-entropy loss: \texttt{loss\ =\ -np.mean(np.log(y\_hat{[}np.arange(n),\ y{]}))}
\item
  Know backprop for a simple 2-layer network --- interviewers at Google/Meta sometimes ask this
\item
  Know how to implement a basic SGD update loop
\end{itemize}

\textbf{ML design questions:}

\begin{itemize}
\tightlist
\item
  Always ask: how much data? Is it labeled? What are the compute constraints? What\textquotesingle s the latency requirement?
\item
  For image tasks: start with pre-trained CNN backbone (EfficientNet/ResNet) + fine-tune
\item
  For text tasks: start with pre-trained BERT/RoBERTa + fine-tune
\item
  Always mention regularization, validation strategy, monitoring
\end{itemize}

\textbf{Conceptual questions:}

\begin{itemize}
\tightlist
\item
  Have first-principles explanations ready, not just "it works because..."
\item
  Be able to derive the math: backprop chain rule, Xavier init formula, cross-entropy gradient
\item
  Connect to practical debugging: "how would you diagnose vanishing gradients?" \(\rightarrow\) monitor gradient norms per layer, check activation distributions
\end{itemize}

\textbf{Behavioral / ML judgment:}

\begin{itemize}
\tightlist
\item
  Know trade-offs: Adam vs SGD, CNN vs Transformer for vision, BatchNorm vs LayerNorm
\item
  Know failure modes: overfitting, underfitting, training instability, data leakage
\item
  Be ready to discuss real systems: recommendation, search ranking, content moderation
\end{itemize}

\textbf{Common FAANG deep learning interview topics by company:}

\begin{itemize}
\tightlist
\item
  \textbf{Google:} Efficiency (EfficientNet, quantization, mobile models), distributed training, TensorFlow/Keras
\item
  \textbf{Meta:} Recommendation systems (DLRM, embeddings), PyTorch internals, ads models
\item
  \textbf{Amazon:} Time-series models (forecasting), on-device ML (Alexa), recommendation
\item
  \textbf{Apple:} On-device ML, CoreML, privacy-preserving training, efficient architectures
\end{itemize}

\section{Revision Cheat Sheet}\label{revision-cheat-sheet}

\subsection{10-Minute Summary}\label{10-minute-summary}

\begin{enumerate}
\def\labelenumi{\arabic{enumi}.}
\item
  \textbf{Perceptron:} \texttt{y\ =\ f(Wx\ +\ b)}. Without \texttt{f}, it\textquotesingle s linear. Non-linearity enables learning complex functions.
\item
  \textbf{MLP:} Stack layers. Each layer = linear transform + activation. Universal approximator with enough neurons.
\item
  \textbf{Forward prop:} Compute \texttt{z\ =\ Wx\ +\ b}, then \texttt{a\ =\ f(z)} for each layer. Store both.
\item
  \textbf{Loss:} MSE for regression. Cross-entropy for classification. Must be differentiable.
\item
  \textbf{Backprop:} Chain rule, layer by layer, right to left. \texttt{dL/dW\ =\ upstream\ *\ local\_derivative\ *\ input}.
\item
  \textbf{Optimizers:} SGD (simple) \(\rightarrow\) +Momentum (accelerate) \(\rightarrow\) RMSprop (adapt LR) \(\rightarrow\) Adam (both). AdamW for transformers.
\item
  \textbf{Vanishing gradients:} Sigmoid saturates \(\rightarrow\) gradient dies. Fix: ReLU, residual connections, BatchNorm, He init.
\item
  \textbf{BatchNorm:} Normalize across batch per feature. Train and inference differ. Good for CNNs.
  \textbf{LayerNorm:} Normalize across features per sample. Same at train/inference. Good for Transformers/RNNs.
\item
  \textbf{Dropout:} Random mask during training, scale to compensate. Disabled at inference. Ensemble effect.
\item
  \textbf{CNNs:} Local, shared filters. Translation invariant. Feature maps = filter responses. Pool for downsampling.
\item
  \textbf{RNNs/LSTMs:} Sequential processing. Vanishing gradient through time. LSTMs: additive cell state = gradient highway. GRU = simpler LSTM.
\item
  \textbf{Transfer learning:} Pre-trained backbone + fine-tune head. Use smaller LR for pre-trained layers.
\end{enumerate}

\subsection{Key Numbers to Remember}\label{key-numbers-to-remember}

\begin{longtable}[]{@{}
  >{\raggedright\arraybackslash}p{(\columnwidth - 2\tabcolsep) * \real{0.3498}}
  >{\raggedright\arraybackslash}p{(\columnwidth - 2\tabcolsep) * \real{0.6502}}@{}}
\toprule\noalign{}
\rowcolor{acc!16}
\begin{minipage}[b]{\linewidth}\raggedright
Quantity
\end{minipage} & \begin{minipage}[b]{\linewidth}\raggedright
Value
\end{minipage} \\
\midrule\noalign{}
\endhead
\bottomrule\noalign{}
\endlastfoot
Sigmoid derivative max & 0.25 (at x=0) \\
\rowcolor{zebra}
Adam defaults & \(\beta\)1=0.9, \(\beta\)2=0.999, \(\varepsilon\)=1e-8 \\
Xavier init variance & 2/(n\_in + n\_out) \\
\rowcolor{zebra}
He init variance & 2/n\_in \\
Typical dropout rate & 0.1 - 0.5 \\
\rowcolor{zebra}
VGG-16 params & \textasciitilde138M \\
ResNet-50 params & \textasciitilde25M \\
\rowcolor{zebra}
LSTM: gates count & 3 (input, forget, output) \\
GRU: gates count & 2 (update, reset) \\
\rowcolor{zebra}
Softmax temperature & 1.0 default; \textless{} 1 = sharper; \textgreater{} 1 = softer \\
\end{longtable}

\subsection{Cheat-Sheet Table: Picking the Right Tool}\label{cheat-sheet-table-picking-the-right-tool}

\begin{longtable}[]{@{}
  >{\raggedright\arraybackslash}p{(\columnwidth - 2\tabcolsep) * \real{0.3723}}
  >{\raggedright\arraybackslash}p{(\columnwidth - 2\tabcolsep) * \real{0.6277}}@{}}
\toprule\noalign{}
\rowcolor{acc!16}
\begin{minipage}[b]{\linewidth}\raggedright
Situation
\end{minipage} & \begin{minipage}[b]{\linewidth}\raggedright
Use
\end{minipage} \\
\midrule\noalign{}
\endhead
\bottomrule\noalign{}
\endlastfoot
Image classification & CNN (ResNet, EfficientNet) \\
\rowcolor{zebra}
Sequence-to-sequence (NLP) & Transformer (BERT, GPT) \\
Time series / sequential & LSTM/GRU or Transformer \\
\rowcolor{zebra}
Small dataset, image task & Pre-trained CNN + freeze + new head \\
Binary output & Sigmoid + BCE loss \\
\rowcolor{zebra}
Multiclass output & Softmax + Cross-Entropy \\
Regression output & Linear + MSE \\
\rowcolor{zebra}
Hidden layers (default) & ReLU \\
Hidden layers (modern arch) & GELU \\
\rowcolor{zebra}
CNN training & BatchNorm \\
Transformer/RNN training & LayerNorm \\
\rowcolor{zebra}
Default optimizer & Adam (or AdamW for transformers) \\
Final accuracy push (vision) & SGD + momentum + cosine schedule \\
\rowcolor{zebra}
Prevent overfitting & Dropout, weight decay, data augmentation, early stopping \\
Training unstable & Gradient clipping, reduce LR, check init \\
\rowcolor{zebra}
Gradients vanishing & ReLU, residual connections, BatchNorm, He init \\
\end{longtable}

\subsection{Most Important Concepts (Interview Priority)}\label{most-important-concepts-interview-priority}

\textbf{P0 --- Must know cold:}

\begin{itemize}
\tightlist
\item
  Forward propagation math
\item
  Backpropagation and chain rule
\item
  Why non-linearity matters
\item
  Vanishing gradient problem and solutions
\item
  Adam optimizer (what each term does)
\item
  BatchNorm vs LayerNorm
\item
  CNN: convolution, filter, stride, pooling, parameter count
\item
  LSTM gating mechanism (additive cell state)
\item
  Dropout mechanics (training vs inference)
\item
  Transfer learning strategy
\end{itemize}

\textbf{P1 --- Should know well:}

\begin{itemize}
\tightlist
\item
  Activation function comparison and trade-offs
\item
  Xavier vs He initialization
\item
  Residual connections (ResNet)
\item
  RNN vanishing gradient (through time)
\item
  Softmax + cross-entropy numerical stability
\item
  Word embeddings intuition
\item
  Overfitting diagnosis and fixes
\end{itemize}

\textbf{P2 --- Good to know for senior roles:}

\begin{itemize}
\tightlist
\item
  Double descent phenomenon
\item
  Adam vs SGD generalization trade-off
\item
  AdamW (decoupled weight decay)
\item
  Gradient checkpointing
\item
  Mixed precision training
\item
  GAN training dynamics
\item
  VAE latent space
\end{itemize}

\emph{This document covers \textasciitilde95\% of deep learning questions encountered in FAANG ML engineer and research scientist interviews. Review the Core Concepts section actively (write code, derive formulas) rather than passively reading.}

\end{document}
